\providecommand{\GL}{\operatorname{GL}}
\providecommand{\SO}{\operatorname{SO}}
\providecommand{\Orth}{\operatorname{O}}
\providecommand{\SU}{\operatorname{SU}}
\providecommand{\U}{\operatorname{U}}
\providecommand{\SL}{\operatorname{SL}}
\providecommand{\SE}{\operatorname{SE}}
\providecommand{\Sym}{\operatorname{Sym}}
\providecommand{\Stab}{\operatorname{Stab}}
\providecommand{\Orb}{\operatorname{Orb}}
\providecommand{\Aut}{\operatorname{Aut}}
\providecommand{\Img}{\operatorname{Im}}
\providecommand{\rank}{\operatorname{rank}}
\providecommand{\tr}{\operatorname{tr}}
\providecommand{\ReLU}{\operatorname{ReLU}}
\providecommand{\Lie}[1]{\mathfrak{#1}}

\chapter{Mathematical Foundations}
\label{ch:foundations}

The premise of geometric deep learning, set out in \Cref{ch:introduction}, is
that the geometric structure of the data should dictate the structure of the
model. Making this premise operational requires two languages developed in
parallel. The first is the language of neural networks: how a network is built
from neurons and layers, what the data it processes really is, and how it
learns. The second is the geometric language: the \emph{symmetries} that a
problem possesses, formalized by groups and their representations, and the
\emph{spaces} on which data lives, formalized by manifolds and homogeneous
spaces. This is the longest chapter of the book precisely because both
foundations must be laid here before the geometric synthesis of \Cref{ch:priors}
and the architectures that follow.

We begin with the artificial neuron and the activation functions that give it
its expressive power. We then ask what the data flowing through a network
actually is, introducing the notion of a signal on a structured domain. We
describe layers and the various ways of viewing a network as a graph. We
formalize learning --- loss functions, gradient descent, backpropagation, the
universal approximation theorem, convergence theory, and modern optimizers ---
and survey the advantages of depth. With the neural foundations in place we turn
to geometry: manifolds and tangent spaces, groups and their linear
representations, and homogeneous spaces, the constructions that unify domains as
different as the sphere and the plane. The treatment is self-contained but
selective; standard references for the geometric material are
\citet{lee2012introduction} for manifolds, \citet{munkres2000topology} for
topology, and \citet{Kirillov2008,Fulton1991} for Lie groups and representation
theory. The heavier measure-theoretic and functional-analytic prerequisites
behind the universal approximation theorem are collected in the appendix
(\Cref{ch:appendix}) and invoked here as needed.

\section{The artificial neuron}
\label{sec:found:neuron}

Artificial neural networks are one of the earliest and most successful
approaches to biologically inspired computation. Their development began in 1943
when Warren McCulloch and Walter Pitts proposed the first formal mathematical
model of an artificial neuron~\citep{McCullochPitts1943}, establishing
conceptual foundations that endure to this day. The central motivation is to
emulate, in simplified form, the information processing of biological neurons:
specialized cells that receive input signals, integrate them through complex
electrochemical processes, and produce output responses that propagate through
interconnected networks. Crucially, biological neurons exhibit synaptic
plasticity, adapting their responses with experience and thereby learning from
examples.

In the nervous system, a neuron receives electrochemical signals from other
cells through dendrites and synaptic connections; these are integrated in the
cell body, and if the weighted sum exceeds a critical depolarization threshold,
the neuron fires an action potential that propagates along the axon. Abstracting
this mechanism, an artificial neuron is a mathematical model that (1) receives
several numerical input signals, (2) weights them through adjustable parameters
(synaptic weights), (3) applies a nonlinear activation function modelling the
firing threshold, and (4) produces an output signal that may feed other neurons.

\begin{figure}[h]
\centering
\begin{tikzpicture}[
    node distance=0.8cm,
    input/.style={circle, draw, minimum size=0.6cm, fill=gdlblue!15},
    weight/.style={rectangle, draw, minimum width=0.6cm, minimum height=0.4cm, fill=gdlgray!15, font=\scriptsize},
    sum/.style={circle, draw, minimum size=0.8cm, fill=gdlyellow!20},
    activation/.style={rectangle, draw, rounded corners, minimum width=0.8cm, minimum height=0.6cm, fill=gdlgreen!20},
    output/.style={circle, draw, minimum size=0.6cm, fill=gdlred!15},
    arrow/.style={->, thick, >=stealth}
]
\node[input] (x1) at (0, 1.5) {$x_1$};
\node[input] (x2) at (0, 0) {$x_2$};
\node[input] (xn) at (0, -1.5) {$x_n$};
\node at (0, -0.75) {$\vdots$};
\node[weight] (w1) at (1.5, 1.5) {$w_1$};
\node[weight] (w2) at (1.5, 0) {$w_2$};
\node[weight] (wn) at (1.5, -1.5) {$w_n$};
\node[sum] (suma) at (3.5, 0) {$\sum$};
\node[below=0.1cm of suma, font=\scriptsize] {$+b$};
\node[activation] (sigma) at (5.5, 0) {$\sigma$};
\node[output] (y) at (7, 0) {$y$};
\draw[arrow] (x1) -- (w1);
\draw[arrow] (x2) -- (w2);
\draw[arrow] (xn) -- (wn);
\draw[arrow] (w1) -- (suma);
\draw[arrow] (w2) -- (suma);
\draw[arrow] (wn) -- (suma);
\draw[arrow] (suma) -- node[above, font=\scriptsize] {$z$} (sigma);
\draw[arrow] (sigma) -- (y);
\node[above=1.5cm of x1, font=\small\bfseries] {Inputs};
\node[above=0.3cm of w1, font=\small\bfseries] {Weights};
\node[above=0.5cm of suma, font=\small\bfseries] {Sum};
\node[above=0.5cm of sigma, font=\small\bfseries] {Activation};
\end{tikzpicture}
\caption[Structure of an artificial neuron]{Structure of an artificial neuron.
The inputs $x_i$ are multiplied by weights $w_i$, summed together with the bias
$b$, and the result $z = \sum_i w_i x_i + b$ is transformed by the activation
function $\sigma$ to produce the output $y = \sigma(z)$.}
\label{fig:found:neuron}
\end{figure}

Real biological neurons are vastly more complex than this abstraction: their
dynamics involve precise temporal patterns, specific chemical neurotransmitters,
time-dependent synaptic plasticity, and intricate three-dimensional dendritic
architectures. Models closer to biology exist, such as \emph{spiking neural
networks}, regarded as the ``third generation'' of neural
networks~\citep{Maass1997spiking}: rather than emitting a continuous value, they
process information through discrete spikes that propagate in time, capturing the
temporal dynamics of biological neurons and offering significant advantages in
energy efficiency on neuromorphic hardware~\citep{Tavanaei2018deep}. The
classical McCulloch--Pitts abstraction has nonetheless proved extraordinarily
fruitful, and it is the model we adopt throughout.

We begin by formalizing the activation function and then the neuron itself.

\begin{definition}[Activation function]
\label{def:found:activation}
An \emph{activation function} is a function $\sigma : \R \to \R$ satisfying:
\begin{itemize}
  \item \textbf{Nonlinearity:} $\sigma$ is not an affine map;
  \item \textbf{Continuity:} $\sigma$ is continuous;
  \item \textbf{Differentiability almost everywhere:} $\sigma$ is differentiable
    except possibly on a set of Lebesgue measure zero.\footnote{A set
    $A \subset \R$ has Lebesgue measure zero if it can be covered by intervals
    whose total length is arbitrarily small; see \Cref{ch:appendix}.}
\end{itemize}
\end{definition}

These properties ensure that $\sigma$ introduces the nonlinearity required for
expressive power while still permitting gradient computation for training by
backpropagation.

\begin{definition}[Artificial neuron]
\label{def:found:neuron}
Let $n \in \mathbb{N}$ and let $\sigma : \R \to \R$ be an activation function.
For each $\theta = (\mathbf{w}, b) \in \R^{n+1}$ with $\mathbf{w} \in \R^n$
(weights) and $b \in \R$ (bias), an \emph{artificial neuron} is the function
$f_\theta : \R^n \to \R$ given by the composition of an affine map with the
activation,
\[
  f_\theta(\mathbf{x}) = \sigma\bigl(T_\theta(\mathbf{x})\bigr)
  = \sigma\bigl(\inner{\mathbf{w}}{\mathbf{x}} + b\bigr),
\]
where $T_\theta(\mathbf{x}) = \inner{\mathbf{w}}{\mathbf{x}} + b$ is the affine
map. A neuron thus defines a \emph{family of functions}
$\{f_\theta : \theta \in \R^{n+1}\}$ parametrized by the weights and bias.
\end{definition}

The nonlinearity of $\sigma$ is essential: a composition of affine maps is again
affine, so without nonlinear activations a deep network with many layers would
collapse to a single linear transformation, losing all expressive power. It is
exactly this nonlinearity that lets neural networks approximate arbitrarily
complex functions, as the universal approximation theorem
(\Cref{thm:found:universal}) makes precise.

\subsection{Desirable properties of activation functions}
\label{sec:found:activation-properties}

\Cref{def:found:activation} states the minimal requirements. In practice it is
desirable that $\sigma$ satisfy additional properties that ease training and
improve network behaviour.

\begin{definition}[Controlled growth]
\label{def:found:bounded}
An activation function $\sigma : \R \to \R$ has \emph{controlled growth} if it
satisfies one of:
\begin{enumerate}
  \item \textbf{Boundedness:} there exists $M > 0$ with $\abs{\sigma(z)} \le M$
    for all $z \in \R$;
  \item \textbf{Linear growth:} there exist $C_1, C_2 > 0$ with
    $\abs{\sigma(z)} \le C_1\abs{z} + C_2$ for all $z \in \R$.
\end{enumerate}
\end{definition}

This property is crucial for numerical stability during training. Boundedness
keeps activations within a finite range regardless of input magnitude,
preventing numerical overflow; linear growth is more permissive, allowing
activations to grow but no faster than the input itself.

\begin{definition}[Monotonicity]
\label{def:found:monotone}
A function $\sigma : \R \to \R$ is \emph{monotone} (increasing) if
$z_1 < z_2$ implies $\sigma(z_1) \le \sigma(z_2)$.
\end{definition}

Monotonicity preserves the relative order of inputs, simplifies the optimization
landscape when combined with a convex loss, and reduces spurious local minima.
Most classical activations are monotone, although monotonicity is not necessary
for good empirical performance.

\begin{definition}[Zero-centred]
\label{def:found:zerocentered}
An activation function $\sigma : \R \to \R$ is \emph{zero-centred} if
$\sigma(0) = 0$ (or $\sigma(0) \approx 0$) and it is odd,
$\sigma(-z) = -\sigma(z)$ for all $z$.
\end{definition}

When activations are zero-centred, gradients have mean close to zero, avoiding
systematic biases in weight updates. When this fails (for instance if
$\sigma(0) \ne 0$), the outputs of each layer carry a systematic bias, so all
gradients with respect to the weights of a layer share the same sign, producing
the ``zigzagging'' behaviour of gradient descent.

\begin{definition}[Non-saturating gradient]
\label{def:found:nonsaturating}
An activation function $\sigma : \R \to \R$ has \emph{non-saturating gradient}
if there exist $\delta > 0$ and a set $S \subseteq \R$ of infinite measure such
that $\abs{\sigma'(z)} \ge \delta$ for all $z \in S$.
\end{definition}

This property is fundamental to avoiding the \emph{vanishing gradient problem},
one of the principal obstacles in training deep networks
(\Cref{sec:found:vanishing}).

\begin{corollary}
\label{cor:found:saturates}
Every bounded activation function with horizontal asymptotes necessarily
saturates.
\end{corollary}

\begin{proof}
If $\lim_{z\to\pm\infty}\sigma(z) = c_{\pm}$ with $c_+, c_-$ finite, then
$\lim_{z\to\pm\infty}\sigma'(z) = 0$. For any $\delta > 0$ there is $M > 0$ with
$\abs{\sigma'(z)} < \delta$ for all $\abs{z} > M$, so the set
$\{z : \abs{\sigma'(z)} \ge \delta\} \subseteq [-M, M]$ has finite measure. By
\Cref{def:found:nonsaturating}, $\sigma$ cannot have a non-saturating gradient,
hence it saturates.
\end{proof}

For large $\abs{z}$, $\abs{\sigma'(z)} \approx 0$, which causes gradients to
become exponentially small in deep layers during backpropagation: the gradient
of the loss with respect to the weights of the first layers is proportional to
$\prod_{\ell=1}^L \sigma'(z^{(\ell)})$, which can vanish numerically. An
activation that keeps $\abs{\sigma'(z)} \ge \delta > 0$ on a set of infinite
measure lets the gradient flow without attenuation across many layers.

\begin{definition}[Computational efficiency]
\label{def:found:efficiency}
An activation function is \emph{computationally efficient} if both $\sigma(z)$
and its derivative $\sigma'(z)$ can be evaluated in $O(1)$ time with basic
arithmetic operations.
\end{definition}

In deep networks with millions of neurons the activation is evaluated billions
of times during training, so its efficiency directly affects total training
time. Functions defined piecewise through $\max$ require only comparisons and
selections, making them extremely efficient; transcendental functions
(exponentials, logarithms) are still $O(1)$ but with a much larger constant.

\begin{definition}[Lipschitz continuity]
\label{def:found:lipschitz}
A function $\sigma : \R \to \R$ is \emph{Lipschitz continuous} if there exists
$L > 0$ (the \emph{Lipschitz constant}) such that
$\abs{\sigma(z_1) - \sigma(z_2)} \le L\abs{z_1 - z_2}$ for all $z_1, z_2 \in \R$.
\end{definition}

Lipschitz continuity provides quantitative control over how input perturbations
propagate to the output: small changes in the weights produce bounded changes in
the activations, and if $\sigma$ is Lipschitz with constant $L$ then
$\abs{\sigma'(z)} \le L$ almost everywhere. Many convergence theorems for
optimization algorithms assume Lipschitz continuity of the objective. For
activations differentiable almost everywhere,
$L = \sup_z \abs{\sigma'(z)}$.

\Cref{tab:found:activations} summarizes the most common activation functions and
their properties; each is described in detail in
\Cref{sec:found:traditional-layers}.

\begin{table}[h]
\centering
\caption{Comparison of common activation functions.}
\label{tab:found:activations}
\begin{tabular}{@{}lccccc@{}}
\toprule
\textbf{Function} & \textbf{Formula} & \textbf{Range} & \textbf{Centred} & \textbf{Non-sat.} & \textbf{Lipschitz} \\
\midrule
Sigmoid & $\frac{1}{1+e^{-z}}$ & $(0,1)$ & No & No & $L=\frac{1}{4}$ \\
\addlinespace
Tanh & $\frac{e^z-e^{-z}}{e^z+e^{-z}}$ & $(-1,1)$ & Yes & No & $L=1$ \\
\addlinespace
ReLU & $\max(0,z)$ & $[0,\infty)$ & No & Yes ($z>0$) & $L=1$ \\
\addlinespace
Leaky ReLU & $\max(\alpha z,z)$ & $\R$ & No & Yes & $L=1$ \\
\addlinespace
Swish & $z \cdot \sigma(z)$ & $\approx(-0.3,\infty)$ & No & Partial & $L\approx 1.1$ \\
\bottomrule
\end{tabular}
\end{table}

\section{Data and its structure}
\label{sec:found:data}

We defined the artificial neuron as a function that processes input data and
produces output data. But we have not yet asked what the data flowing through a
network really is --- what the space $\R^n$ actually represents. When we
represent a $28 \times 28$ image as a vector in $\R^{784}$, we are not dealing
with $784$ independent values: each component corresponds to a specific pixel on
a two-dimensional grid, and spatially neighbouring pixels are naturally related.
The grid structure is not arbitrary but intrinsic to the image. This section
formalizes that distinction through two equivalent conceptual frameworks: the
\emph{domain--structure--signal} triple $(\Omega, \mathcal{S}, f)$ and the
\emph{geometry--attributes} pair.

\subsection{Signals on structured domains}
\label{sec:found:signals}

The first framework rests on three components that cleanly separate the notions
of position, relation, and value. This tripartite decomposition underlies the
unifying paradigm developed in \Cref{ch:priors} and the chapters on the five
themes.

\begin{definition}[Data domain]
\label{def:found:domain}
A \emph{data domain} is a set $\Omega$ whose elements represent the positions
where data is observed. It may be discrete and finite
($\Omega = \{p_1, \ldots, p_n\}$, e.g.\ graph vertices, image pixels, vector
indices), discrete and infinite ($\Omega = \Z^d$, an infinite grid), or
continuous ($\Omega \subseteq \R^d$ or a smooth manifold $\mathcal{M}$).
\end{definition}

\begin{definition}[Domain structure]
\label{def:found:structure}
The \emph{structure} $\mathcal{S}$ on a domain $\Omega$ specifies the intrinsic
relations between its elements. It may be \emph{topological} (a topology defining
open sets and neighbourhoods), \emph{metric} (a distance
$d : \Omega \times \Omega \to \R_{\ge 0}$), \emph{algebraic} (a group $\group$
acting on $\Omega$), \emph{combinatorial} (a graph $(V,E)$ of explicit
connections), or \emph{geometric} (a smooth manifold with atlas, Riemannian
metric, connection). The pair $(\Omega, \mathcal{S})$ is a \emph{structured
domain}.
\end{definition}

\begin{definition}[Signal]
\label{def:found:signal}
A \emph{signal} on a structured domain $(\Omega, \mathcal{S})$ is a function
$f : \Omega \to \mathcal{C}$ assigning to each point $p \in \Omega$ a feature
vector $f(p) \in \mathcal{C}$, where $\mathcal{C}$ is a vector space called the
\emph{channel space}, typically $\mathcal{C} = \R^d$ with $d$ the number of
channels. When $\Omega = \{p_1, \ldots, p_n\}$ is finite and
$\mathcal{C} = \R^d$, the signal is represented by the matrix
$F = (f(p_1)^\top; \ldots; f(p_n)^\top) \in \R^{n \times d}$, each row holding
the features observed at position $p_i$.
\end{definition}

\begin{remark}
\label{rem:found:channels}
The choice $\mathcal{C} = \R^d$ suffices for most of this chapter. In later
chapters the channel space acquires richer structure: in
\Cref{ch:groups}, $\mathcal{C}$ may be a representation space of a symmetry group
$\group$ (e.g.\ spherical harmonics for $\SO(3)$), and for gauge-equivariant
networks $\mathcal{C}$ is the fibre of an associated vector bundle over the base
manifold.
\end{remark}

\begin{definition}[Signal space]
\label{def:found:signal-space}
Given a structured domain $(\Omega, \mathcal{S})$ and a channel space
$\mathcal{C}$, the \emph{signal space} is
$\mathcal{F}(\Omega, \mathcal{C}) = \{f : \Omega \to \mathcal{C}\}$. When
$\Omega$ is finite with $\abs{\Omega} = n$ and $\mathcal{C} = \R^d$, it is
isomorphic to the matrix space $\R^{n \times d}$. When $\Omega$ is continuous,
$\mathcal{F}(\Omega, \mathcal{C})$ restricts to a function space with appropriate
regularity (for example $L^2(\Omega, \mathcal{C})$ or
$C^k(\Omega, \mathcal{C})$).
\end{definition}

A neural network can then be seen as an operator between signal spaces: a layer
transforms input signals into output signals, possibly over different domains
and channel spaces, $T : \mathcal{F}(\Omega, \mathcal{C}) \to
\mathcal{F}(\Omega', \mathcal{C}')$. This functional viewpoint is central to the
convolutional networks of \Cref{ch:grids}, where convolution is an equivariant
linear operator between signal spaces. In summary: the domain $\Omega$ defines
\emph{where} the data is, the structure $\mathcal{S}$ defines \emph{how the
positions relate}, and the signal $f$ defines the \emph{values observed} at each
position. A vector $\mathbf{x} \in \R^n$ is therefore not merely a collection of
$n$ numbers but a signal on a domain with a specific structure.

\begin{example}[Audio signal]
\label{ex:found:audio}
An audio signal has domain $\Omega = [0,T] \subset \R$ (a continuous time
interval), structure $\mathcal{S}$ given by the total order of $\R$ with its
standard topology, and signal $f(t) \in \R$ (sound amplitude at instant $t$).
The order is causal: $t_1 < t_2$ has meaning, and nearby instants are typically
correlated.
\end{example}

\begin{example}[Digital image on a grid]
\label{ex:found:image}
A digital image has domain $\Omega = \{1,\ldots,H\} \times \{1,\ldots,W\}$,
structure given by a rectangular 2D grid with local (4- or 8-connected)
neighbourhoods, and signal $f(i,j) \in \R^3$ (RGB intensities). Permuting the
pixel order when vectorizing changes the vector representation but not the image
as a signal on the grid: the grid structure is intrinsic to the geometric
object, not to its particular vectorization.
\end{example}

\begin{example}[Social network]
\label{ex:found:graph}
A social network has domain $\Omega = V = \{u_1, \ldots, u_n\}$ (users),
structure $\mathcal{S} = (V, E)$ with $E \subseteq V \times V$ representing
friendships, and signal $f(u) \in \R^d$ (the profile of user $u$). There is no
canonical ordering of users: a user's ``position'' is its abstract identity in
the graph, and two users are neighbours if and only if $(u_i, u_j) \in E$.
\end{example}

\begin{example}[Molecule]
\label{ex:found:molecule}
A molecule combines discrete connectivity with continuous geometry: domain
$\Omega = \{a_1, \ldots, a_n\}$ (atoms), structure given by the chemical-bond
graph together with positions in $\R^3$ and the symmetries of the Euclidean
group $E(3)$, and signal $f(a) \in \R^d$ (atomic type, partial charge, mass,
hybridization). Physicochemical properties are invariant under $E(3)$:
translations, rotations, and reflections do not change the molecule.
\end{example}

\begin{example}[3D point cloud and surface]
\label{ex:found:pointcloud}
A 3D point cloud has domain $\Omega = \{p_1, \ldots, p_n\} \subset \R^3$,
structure the Euclidean metric $d(p_i, p_j) = \norm{p_i - p_j}_2$, and signal
$f(p) \in \R^d$ (surface normals, curvatures, colours, LIDAR intensity). There is
no natural order: only relative distances matter. A 3D surface modelled as a
Riemannian manifold $\mathcal{M}$ embedded in $\R^3$ carries the induced metric
$g$, and its geodesic distances and curvature are intrinsic, independent of the
particular embedding.
\end{example}

\begin{remark}[Permutation invariance]
\label{rem:found:permutation}
The examples reveal a fundamental distinction. In some domains (audio, text,
image) there is a canonical or partial ordering of positions, so the matrix
representation $F \in \R^{n \times d}$ is essentially unique. In others (graphs,
point clouds) there is no such order: if $P \in \{0,1\}^{n \times n}$ is a
permutation matrix, then $F$ and $PF$ represent the same signal on the same
domain, merely with positions enumerated in a different order. A network
processing data without a canonical order must therefore produce results that are
invariant (for classification) or equivariant (for per-node prediction) under
input permutations --- one of the central ideas of geometric deep learning,
formalized in \Cref{ch:priors,ch:graphs}.
\end{remark}

\subsection{Geometry and attributes}
\label{sec:found:geometry-attributes}

The previous framework separates domain, structure, and signal into three
components. It is sometimes useful to unify the domain and its structure into a
single entity: the \emph{geometry} of the problem.

\begin{definition}[Geometry and attributes]
\label{def:found:geometry}
The \emph{geometry} of a data problem is the pair $\mathcal{D} = (\Omega,
\mathcal{S})$ combining the domain with its structure into a single entity. The
\emph{attributes} (or features) are the signal $f : \Omega \to \mathcal{C}$
assigning observed values to each position.
\end{definition}

The two frameworks are mathematically equivalent, with the correspondence
$(\Omega, \mathcal{S}, f) \leftrightarrow (\mathcal{D}, f)$. The
geometry/attributes view emphasizes that the geometry $\mathcal{D}$ is typically
fixed or intrinsic to each datum (it defines the ``shape'' of the data space),
while the attributes $f$ are variable: different data instances share the same
geometry but have different attributes. Operations should respect the geometry
while transforming the attributes.

\subsection{Neural processing of signals}
\label{sec:found:neuron-as-operator}

An artificial neuron $f_\theta : \R^n \to \R$ can be reinterpreted as an operator
on signals over a discrete domain: input domain $\Omega = \{1, \ldots, n\}$ (a
finite set of indices with no particular geometric meaning), no significant
structure, input signal $x : \Omega \to \R$ with $x(i) = x_i$, and processing
$y = \sigma(\inner{\mathbf{w}}{\mathbf{x}} + b)$. The linear part
$\inner{\mathbf{w}}{\mathbf{x}} = \sum_i w_i x_i$ combines all components
globally. There is no notion of \emph{locality} (all positions contribute
equally), no \emph{structured neighbourhood} (the connection $w_1 x_1$ is
structurally indistinguishable from $w_{100} x_{100}$), and no \emph{geometric
invariance} (permuting the indices generally changes the output, even when the
permutation is semantically irrelevant). Traditional networks focus on
transforming the signal between layers, leaving the geometry in the background;
geometric deep learning makes the geometry of the data a central pillar, as
developed in \Cref{ch:priors}.

\section{Layers of a neural network}
\label{sec:found:layers}

We defined an artificial neuron as a function with a single scalar output. This
is not arbitrary: biological neurons emit a single output signal (firing rate),
and a $\R^n \to \R$ function is conceptually simpler than a $\R^n \to \R^m$ one,
easing the analysis of an individual unit. A vector-valued neuron adds no
expressive power: a neuron with $m$ outputs is functionally equivalent to $m$
scalar neurons operating in parallel. We therefore take the scalar neuron as the
primitive and build richer architectures by composition and parallelization.

\begin{definition}[Neural layer]
\label{def:found:layer}
A \emph{neural layer} is a set of $m$ neurons sharing the same input domain and
operating in parallel. Given an input signal $\mathbf{x} \in \R^n$, the layer
produces $\mathbf{y} \in \R^m$ via
\[
  L(\mathbf{x}) = \sigma(\mathbf{W}\mathbf{x} + \mathbf{b}),
\]
where $\mathbf{W} \in \R^{m \times n}$ is the \emph{weight matrix} (whose $i$-th
row holds the weights of neuron $i$), $\mathbf{b} \in \R^m$ is the \emph{bias
vector}, and $\sigma$ is applied componentwise. Each neuron computes
$y_i = \sigma(\mathbf{w}_i^\top \mathbf{x} + b_i)$.
\end{definition}

The matrix form is fully general: if a neuron processes only a subset of input
positions, this is expressed by zero coefficients $W_{ij} = 0$. Architectures
with structured connectivity, such as convolutional networks, exploit this
sparsity. In many architectures the weights are \emph{shared} between neurons:
several rows of $\mathbf{W}$ hold identical values, drastically reducing the
number of free parameters.

\subsection{Composition of layers}
\label{sec:found:composition}

A single neuron computes $y = \sigma(\sum_i w_i x_i + b)$, whose linear part
defines a hyperplane in $\R^n$; a single neuron can only separate the input space
by one nonlinearly transformed hyperplane, limiting its expressive power. By
connecting many neurons in layers we obtain a \emph{functional composition} of
nonlinear transformations. Writing $L_i : \R^{n_{i-1}} \to \R^{n_i}$ for the
$i$-th layer, a $k$-layer network computes
\[
  \mathcal{N} = L_k \circ L_{k-1} \circ \cdots \circ L_1 :
  \R^{n_0} \to \R^{n_k}.
\]
This composition yields \emph{hierarchical representations} (early layers detect
simple patterns such as edges and textures; later layers combine them into
higher-level structures) and \emph{complex partitioning} of the space (multiple
neurons generate multiple hyperplanes, partitioning $\R^n$ into arbitrarily
complex polyhedral regions). The first layer is the input layer, the last is the
output layer, and the intermediate layers are the \emph{hidden} layers.

\begin{figure}[h]
\centering
\begin{tikzpicture}[
    x=2.1cm, y=0.85cm,
    neuron/.style={circle, minimum size=5.5mm, line width=0.6pt},
    input/.style={neuron, draw=gdlblue!70, fill=gdlblue!15},
    hidden1/.style={neuron, draw=gdlgreen!60!black, fill=gdlgreen!15},
    hidden3/.style={neuron, draw=gdlpurple!70, fill=gdlpurple!18},
    output/.style={neuron, draw=gdlorange!70, fill=gdlorange!18},
    conn/.style={->, >=stealth, draw=gdlgray!35, line width=0.2pt},
    lbl/.style={font=\small\bfseries}
]
\node[input] (I1) at (0, 2) {};
\node[input] (I2) at (0, 1) {};
\node[input] (I3) at (0, 0) {};
\node[left=1mm of I1, font=\small] {$x_1$};
\node[left=1mm of I3, font=\small] {$x_3$};
\node[lbl] at (0, 3.5) {Input};
\foreach \j in {1,...,4} \node[hidden1] (H1-\j) at (1, 3-\j) {};
\node[lbl] at (1, 3.5) {Hidden 1};
\foreach \j in {1,...,4} \node[hidden1] (H2-\j) at (2, 3-\j) {};
\node[lbl] at (2, 3.5) {Hidden 2};
\foreach \j in {1,...,4} \node[hidden3] (H3-\j) at (3, 3-\j) {};
\node[lbl] at (3, 3.5) {Hidden 3};
\node[output] (O1) at (4, 1.5) {};
\node[output] (O2) at (4, 0.5) {};
\node[right=1mm of O1, font=\small] {$y_1$};
\node[right=1mm of O2, font=\small] {$y_2$};
\node[lbl] at (4, 3.5) {Output};
\foreach \i in {1,2,3} \foreach \j in {1,...,4} \draw[conn] (I\i) -- (H1-\j);
\foreach \i in {1,...,4} \foreach \j in {1,...,4} \draw[conn] (H1-\i) -- (H2-\j);
\foreach \i in {1,...,4} \foreach \j in {1,...,4} \draw[conn] (H2-\i) -- (H3-\j);
\foreach \i in {1,...,4} \foreach \j in {1,2} \draw[conn] (H3-\i) -- (O\j);
\draw[decorate, decoration={brace, amplitude=4pt, raise=2pt}, gdlgray!60, thick]
    (1, 3.7) -- (3, 3.7) node[midway, above=5pt, font=\small\itshape] {Hidden layers};
\end{tikzpicture}
\caption[Structure of a layered neural network]{A feedforward neural network
organized in layers. Nodes are neurons coloured by layer type; arrows are
weighted connections $w_{ij}^{(\ell)}$.}
\label{fig:found:network}
\end{figure}

\paragraph{The XOR problem.}
The exclusive-or function $\mathrm{XOR}(x_1, x_2)$ returns $1$ when exactly one
input is $1$ and $0$ otherwise. It is \emph{not linearly separable}: no single
hyperplane separates the classes $\{(0,0),(1,1)\}$ from $\{(0,1),(1,0)\}$ in
$\R^2$. A single neuron computes one hyperplane and cannot solve it; many layers
of one neuron each also fail, because a scalar map $\R^2 \to \R$ collapses the
two-dimensional information immediately. A network with two hidden neurons and
one output neuron suffices: the hidden layer creates two hyperplanes that
partition $\R^2$, and the output neuron combines them to produce the XOR
boundary. Thus neither width alone nor depth alone is enough: we need at least
two neurons in some layer \emph{and} at least two layers.

\subsection{Width versus depth}
\label{sec:found:width-depth}

\begin{definition}[Width and depth]
\label{def:found:width-depth}
The \emph{width} of a layer $L_i : \R^{n_{i-1}} \to \R^{n_i}$ is the dimension
$n_i$ of its output space, the number of neurons it contains. The \emph{depth} of
a network $\mathcal{N} = L_k \circ \cdots \circ L_1$ is the number $k$ of layers
in the composition.
\end{definition}

Width provides \emph{feature diversity}: each neuron captures a different aspect
of the input. Depth provides \emph{hierarchical abstraction}: each layer performs
an intermediate transformation, building a hierarchy of increasingly abstract
representations,
$x \xrightarrow{L_1} h^{(1)} \xrightarrow{L_2} \cdots \xrightarrow{L_k} y$. Early
layers detect low-level features (edges, textures), intermediate layers combine
them into more complex patterns, and final layers represent task-specific
high-level concepts. As the universal approximation theorem
(\Cref{thm:found:universal}) shows, a single sufficiently wide hidden layer can
approximate any continuous function on a compact set, but this existential
guarantee says nothing about the efficiency of the approximation; depth provides
exponential advantages in expressive power (\Cref{sec:found:depth}).

\subsection{Types of layers}
\label{sec:found:layer-types}

Layers may transform the signal, the domain, or both. This distinction is
fundamental to the difference between traditional networks and geometric deep
learning: the former typically ignore or destroy the original geometry, while the
latter preserve and exploit it.

\begin{definition}[Pure signal transformation]
\label{def:found:signal-transform}
A layer $L$ is a \emph{pure signal transformation} if it preserves the structure
of the input domain: if the input is $f : \Omega \to \R^{c_{\mathrm{in}}}$, the
output is $g : \Omega \to \R^{c_{\mathrm{out}}}$ with $\Omega$ unchanged.
Typically $g(x) = \phi(f(x))$ for some
$\phi : \R^{c_{\mathrm{in}}} \to \R^{c_{\mathrm{out}}}$ applied pointwise; such
layers preserve the tensor shape and require no neighbourhood information.
\end{definition}

\begin{definition}[Pure geometric transformation]
\label{def:found:geometry-transform}
A layer $L$ is a \emph{pure geometric transformation} if it modifies the
geometric structure of the domain without nontrivially processing the signal
values: $L(f) = g : (\Omega', \mathcal{G}') \to \R^c$ with
$\mathcal{G}' \ne \mathcal{G}$, the values being preserved or reindexed by a
correspondence $\pi : \Omega' \to \Omega$.
\end{definition}

\begin{definition}[Coupled signal--geometry transformation]
\label{def:found:coupled-transform}
A layer $L$ is a \emph{coupled transformation} if it simultaneously and
intrinsically modifies both signal and domain: $L$ maps
$f : \Omega \to \R^{c_{\mathrm{in}}}$ to $g : \Omega' \to \R^{c_{\mathrm{out}}}$
with $\Omega' \ne \Omega$ and $g$ not expressible as a reindexing of $f$.
\end{definition}

Traditional networks (multilayer perceptrons, dense networks) rely mainly on
layers that destroy or ignore the original geometry, typically flattening it into
a vector and losing all information about spatial neighbourhoods. Geometric deep
learning, by contrast, systematically preserves and exploits the underlying
geometry, as developed throughout the book.

\subsection{Traditional layers}
\label{sec:found:traditional-layers}

We briefly present some layers used traditionally in deep learning.

\paragraph{Activation functions as layers.}
Although the activation is intrinsic to the neuron, it is often implemented as a
separate layer applying $\sigma$ componentwise; these are pure signal
transformations. The \emph{sigmoid}
$\sigma(z) = 1/(1+e^{-z})$ has range $(0,1)$, is bounded
(\Cref{def:found:bounded}), monotone (\Cref{def:found:monotone}), and Lipschitz
with $L = 1/4$, but is not zero-centred (\Cref{def:found:zerocentered}) and
saturates at both ends. The \emph{hyperbolic tangent}
$\tanh(z) = (e^z - e^{-z})/(e^z + e^{-z})$ has range $(-1,1)$, is zero-centred
and Lipschitz with $L = 1$, but, being bounded with horizontal asymptotes, also
saturates --- causing the vanishing gradient problem of \Cref{sec:found:vanishing}.
The \emph{rectified linear unit} $\ReLU(z) = \max(0, z)$ has
$\ReLU'(z) = 1$ for $z > 0$, satisfying \Cref{def:found:nonsaturating} on the
positive half-line; it has linear growth, is monotone and Lipschitz with $L=1$,
and is extremely efficient, but is not zero-centred and its zero gradient for
$z \le 0$ can cause ``dead neurons.'' The \emph{leaky ReLU}
$\mathrm{LeakyReLU}(z) = \max(\alpha z, z)$ with small $\alpha > 0$ keeps
$\abs{\sigma'(z)} \ge \alpha > 0$ for all $z \ne 0$, eliminating both saturation
and dead neurons.

\begin{definition}[Softmax]
\label{def:found:softmax}
The \emph{softmax} function $\sigma : \R^n \to (0,1)^n$ transforms a real vector
into a probability distribution: for $\mathbf{z} \in \R^n$,
\[
  \sigma(\mathbf{z})_i = \frac{e^{z_i}}{\sum_{j=1}^n e^{z_j}}, \qquad
  i = 1, \ldots, n.
\]
\end{definition}

\begin{proposition}[Properties of softmax]
\label{prop:found:softmax}
The softmax satisfies: (1) \emph{normalization}, $\sum_i \sigma(\mathbf{z})_i =
1$; (2) \emph{positivity}, $\sigma(\mathbf{z})_i > 0$; (3) \emph{translation
invariance}, $\sigma(\mathbf{z} + c\mathbf{1}) = \sigma(\mathbf{z})$; (4) the
Jacobian is $\partial \sigma_i / \partial z_j = \sigma_i(\delta_{ij} -
\sigma_j)$; and (5) \emph{permutation equivariance}, $\sigma(P\mathbf{z}) =
P\sigma(\mathbf{z})$ for every permutation matrix $P$.
\end{proposition}

\begin{proof}
Properties (1) and (2) are immediate. For (3), $e^{z_i + c}/\sum_j e^{z_j + c} =
e^{z_i}/\sum_j e^{z_j}$. For (4), the quotient rule gives
$\partial \sigma_i/\partial z_j = (\delta_{ij}e^{z_i}\sum_k e^{z_k} -
e^{z_i}e^{z_j})/(\sum_k e^{z_k})^2 = \sigma_i(\delta_{ij} - \sigma_j)$.
Property (5) follows because the denominator is invariant under reindexing.
\end{proof}

\paragraph{Normalization.}
Normalization layers standardize the distribution of activations, addressing
internal covariate shift and gradient instability; they are pure signal
transformations. A normalization applies
$\mathrm{Norm}(\mathbf{x}) = \gamma \odot (\mathbf{x} - \boldsymbol{\mu})/
\boldsymbol{\sigma} + \boldsymbol{\beta}$, where $\boldsymbol{\mu},
\boldsymbol{\sigma}$ are centring and scaling statistics and $\gamma,
\boldsymbol{\beta}$ are learnable. \emph{Batch
normalization}~\citep{ioffe2015batch} computes statistics over the mini-batch;
\emph{layer normalization}~\citep{ba2016layer} over the features of each sample;
\emph{group normalization}~\citep{wu2018group} over groups of channels; and
\emph{instance normalization}~\citep{ulyanov2016instance} per channel and per
sample.

\paragraph{Dropout.}
Dropout~\citep{srivastava2014dropout} is a regularization technique that, during
training, randomly deactivates neurons with probability $p_{\mathrm{drop}}$,
setting their outputs to zero:
$\mathbf{y} = \frac{1}{1-p_{\mathrm{drop}}}\,\mathbf{r} \odot \mathbf{x}$ with
$\mathbf{r} \sim \mathrm{Bernoulli}(1-p_{\mathrm{drop}})$. It reduces
co-adaptation between neurons and improves generalization. It is a pure signal
transformation.

\paragraph{Dense layers.}
A \emph{dense} (fully connected) layer is precisely a neural layer
(\Cref{def:found:layer}) with no structural restriction on $\mathbf{W}$: every
output neuron receives signal from every input position. It performs a linear
transformation, a translation, and a nonlinear deformation, and by imposing no
sparse connectivity or weight sharing it destroys the original geometry of the
input through an affine projection. It is a coupled signal--geometry
transformation.

\paragraph{Embedding layers.}
An \emph{embedding} maps discrete indices to dense vectors:
$\mathrm{Embedding} : \{1, \ldots, V_{\mathrm{vocab}}\} \to \R^{d_{\mathrm{emb}}}$,
implemented by a lookup matrix
$E \in \R^{V_{\mathrm{vocab}} \times d_{\mathrm{emb}}}$. Embeddings transform a
discrete, disconnected space into a continuous manifold where Euclidean geometry
captures semantic relations; they are coupled transformations (discrete
$\to$ continuous).

\paragraph{Structure transformations.}
\emph{Reshape} reorganizes a tensor while preserving the total number of
elements; \emph{flatten} is the special case mapping
$\R^{d_1 \times \cdots \times d_n} \to \R^{\prod_i d_i}$, often used to bridge
convolutional and dense layers. \emph{Pooling} reduces spatial dimension by
aggregating local windows with a function $f$, providing limited translational
invariance: \emph{max pooling} selects the maximum, \emph{average pooling} the
arithmetic mean.

\paragraph{Convolutions.}
Convolutional layers, studied in detail in \Cref{ch:grids}, are used mainly for
grid data such as images or audio. For an input $f : \Z^d \to \R$ and a kernel
$g : \Z^d \to \R$, the \emph{discrete convolution} is
$(f * g)[n] = \sum_{m \in \Z^d} f[m]\, g[n - m]$. In traditional networks
convolutions are introduced chiefly as a computational optimization exploiting
the regular grid; \Cref{ch:grids} develops the geometric meaning.

\paragraph{Recurrent layers.}
A \emph{recurrent neural network} has at least one cycle in its connectivity
graph, with hidden state evolving as
$\mathbf{h}_t = f(\mathbf{h}_{t-1}, \mathbf{x}_t; \boldsymbol{\theta})$. The
simplest form uses
$\mathbf{h}_t = \tanh(W_h \mathbf{h}_{t-1} + W_x \mathbf{x}_t + \mathbf{b}_h)$,
$\mathbf{y}_t = W_y \mathbf{h}_t + \mathbf{b}_y$, with the same parameters applied
at every time step. This gives temporal generalization and parameter efficiency
but suffers from limited memory and vanishing gradients, and requires
backpropagation through time~\citep{werbos1990bptt}.

\paragraph{Long short-term memory.}
\emph{Long short-term memory} (LSTM) cells~\citep{hochreiter1997long} address the
vanishing gradient by maintaining a cell state $\mathbf{C}_t$ and a hidden state
$\mathbf{h}_t$ controlled by three gates --- forget
$\mathbf{f}_t = \sigma(W_f [\mathbf{h}_{t-1}, \mathbf{x}_t] + \mathbf{b}_f)$,
input $\mathbf{i}_t$, and output $\mathbf{o}_t$ --- with updates
\[
  \tilde{\mathbf{C}}_t = \tanh(W_C [\mathbf{h}_{t-1}, \mathbf{x}_t] + \mathbf{b}_C),
  \quad \mathbf{C}_t = \mathbf{f}_t \odot \mathbf{C}_{t-1} + \mathbf{i}_t \odot \tilde{\mathbf{C}}_t,
  \quad \mathbf{h}_t = \mathbf{o}_t \odot \tanh(\mathbf{C}_t).
\]
The forget gate can stay near $1$ for important information, creating gradient
paths that persist over long horizons. The \emph{gated recurrent unit}
(GRU)~\citep{cho2014learning} is a simplification combining the forget and input
gates, with comparable performance~\citep{chung2014empirical}.

\paragraph{Residual connections.}
A \emph{residual connection}~\citep{he2016deep} adds the input of a block of
layers $\mathcal{F}$ directly to its output, $\mathbf{y} = \mathcal{F}(\mathbf{x})
+ \mathbf{x}$. The skip connection lets the gradient flow directly during
backpropagation, mitigating the vanishing gradient in very deep networks and
enabling the training of networks far deeper than previously possible.

\section{Neural networks as graphs}
\label{sec:found:nets-as-graphs}

We have described layers and their composition algebraically. There is an equally
important complementary view: representing neural networks as \emph{graphs}, with
nodes as computational units and edges as the flow of information. This view
gives a precise language for architecture, distinguishes families of networks by
topological properties, and underlies the automatic differentiation of modern
frameworks.

\begin{definition}[Graph and directed graph]
\label{def:found:graph}
A \emph{graph} is a pair $\mathcal{G} = (V, E)$ of a vertex set $V$ and an edge
set $E$. In a \emph{directed graph} each edge $(u,v) \in \vec{E}$ has a direction
from $u$ to $v$, and $(u,v) \ne (v,u)$; an \emph{undirected} graph has symmetric
edges. A \emph{path} of length $k$ is a sequence $v_0, \ldots, v_k$ with
$(v_i, v_{i+1}) \in \vec{E}$; a \emph{cycle} is a path with $v_0 = v_k$,
$k \ge 1$, and distinct intermediate vertices. A \emph{directed acyclic graph}
(DAG) is a directed graph with no cycles.
\end{definition}

In a neural network the directed edges represent the flow of information: the
output of neuron $u$ becomes an input to neuron $v$. A cycle represents feedback
or memory.

\begin{theorem}[Topological ordering]
\label{thm:found:topological}
Every directed acyclic graph $\mathcal{G} = (V, \vec{E})$ admits at least one
\emph{topological ordering}: a numbering $v_1, \ldots, v_{\abs{V}}$ such that
$(v_i, v_j) \in \vec{E}$ implies $i < j$.
\end{theorem}

\begin{proof}
By induction on $n = \abs{V}$. For $n = 1$ the ordering is trivial. Assume the
result for $n$ vertices and let $\mathcal{G}$ have $n+1$. We claim $\mathcal{G}$
has a vertex $v$ of in-degree zero: otherwise every vertex would have an incoming
edge, and walking backwards indefinitely would, by the pigeonhole principle,
produce a cycle, contradicting acyclicity. Place such a $v$ first; the reduced
graph on $V \setminus \{v\}$ is still acyclic and has $n$ vertices, so by
induction it has a topological ordering. Prepending $v$ gives one for
$\mathcal{G}$.
\end{proof}

A topological ordering guarantees that the outputs of all neurons can be computed
sequentially without solving implicit systems, making feedforward networks
computationally efficient.

\paragraph{Feedforward networks as DAGs; recurrent networks as cyclic graphs.}
Feedforward networks are exactly those whose neuron graph is a DAG. Acyclicity
allows deterministic one-pass computation following a topological order, parallel
computation of neurons at the same depth, and efficient backpropagation in a
single reverse pass. Recurrent networks have directed cycles representing memory,
yielding equations $h_t = f(h_{t-1}, x_t; \theta)$. Modern architectures tend to
favour DAGs for parallelizability and numerical stability; attention and
transformers (\Cref{sec:found:modern}) capture long-range dependencies while
preserving the DAG structure.

\paragraph{Layer graphs.}
At a higher level of abstraction, a network may be viewed as a \emph{layer
graph} $\mathcal{G}_{\mathcal{N}} = (V_L, E_L)$, with each node a complete layer
and an edge $(L_i, L_j)$ when the output of $L_i$ is part of the input of $L_j$.
A classical feedforward network is the simplest possible DAG, a directed path
$L_1 \to L_2 \to \cdots \to L_k$. Modern architectures use nonlinear DAGs:
residual connections~\citep{he2016deep} add skip edges
$h^{(\ell+1)} = h^{(\ell)} + F(h^{(\ell)}; \theta_\ell)$, and Inception-style
networks~\citep{szegedy2015going} process the input through parallel paths that
are later concatenated. The neuron graph is the microscopic description, the
layer graph the macroscopic one; a projection $\pi : V_N \to V_L$ relates them.

\section{Learning in neural networks}
\label{sec:found:learning}

Training a neural network is an iterative process of optimizing its parameters
$\theta$ to minimize an error measured against examples of the expected output.
With the error at each step, an optimization algorithm --- gradient descent or
one of its variants such as stochastic gradient descent --- updates $\theta$.

\subsection{Loss functions}
\label{sec:found:loss}

\begin{definition}[Loss function]
\label{def:found:loss}
A \emph{loss function} is a map $\mathcal{L} : \mathcal{Y} \times \mathcal{Y} \to
\R_{\ge 0}$ measuring the discrepancy between a prediction $\hat{y}$ and a true
value $y$, satisfying: (1) \emph{non-negativity}, $\mathcal{L}(y, \hat{y}) \ge
0$; (2) \emph{consistency}, $\mathcal{L}(y, y) = 0$; and (3)
\emph{differentiability} with respect to $\hat{y}$, to permit gradient-based
optimization.
\end{definition}

Some losses satisfy the stronger condition $\mathcal{L}(y, \hat{y}) = 0
\Leftrightarrow y = \hat{y}$ (e.g.\ mean squared error), while others reach zero
on a larger region. For \emph{regression}, the \emph{mean squared error}
$\mathcal{L}_{\mathrm{MSE}} = \frac{1}{n}\sum_i (y_i - \hat{y}_i)^2$ penalizes
large errors strongly and is sensitive to outliers; the \emph{mean absolute
error} $\frac{1}{n}\sum_i \abs{y_i - \hat{y}_i}$ is more robust but not
differentiable at $\hat{y} = y$; the \emph{Huber loss} is quadratic for small
errors and linear for large ones, combining the advantages of both. For
\emph{classification}, the \emph{cross-entropy}
$\mathcal{L}(y, \hat{y}) = -\sum_{i=1}^K y_i \log \hat{y}_i$ measures the
divergence between the true and estimated class distributions (with softmax
producing the probabilities), specializing to the \emph{binary log loss} for
$K = 2$; the \emph{hinge loss} $\max(0, 1 - y\hat{y})$ induces maximum-margin
solutions. For comparing \emph{probability distributions}, the
\emph{Kullback--Leibler divergence}
$D_{\mathrm{KL}}(P \| Q) = \sum_x P(x) \log(P(x)/Q(x))$ is a special case of the
\emph{$f$-divergence} $D_f(P \| Q) = \sum_x Q(x) f(P(x)/Q(x))$ for a convex $f$
with $f(1) = 0$, which also encompasses Jensen--Shannon, total variation, and
$\chi^2$. For \emph{representation learning}, the \emph{contrastive
loss}~\citep{hadsell2006dimensionality} and the \emph{triplet
loss}~\citep{schroff2015facenet} operate on pairs and triplets of embeddings,
shaping the geometry so that similar items are close and dissimilar ones are
far.

\begin{definition}[Metric]
\label{def:found:metric}
A \emph{metric} on a set $X$ is a map $d : X \times X \to \R_{\ge 0}$ satisfying,
for all $x, y, z$: (1) non-negativity; (2) identity of indiscernibles,
$d(x,y) = 0 \Leftrightarrow x = y$; (3) symmetry, $d(x,y) = d(y,x)$; and (4) the
triangle inequality $d(x,z) \le d(x,y) + d(y,z)$.
\end{definition}

Many losses are not metrics: cross-entropy is not symmetric, and the
Kullback--Leibler divergence satisfies neither symmetry nor the triangle
inequality. What matters for learning is not that $\mathcal{L}$ be a metric, but
that it provide informative gradients, penalize large errors, and align with the
task.

\subsection{Gradient descent}
\label{sec:found:gradient-descent}

Supervised training by backpropagation proceeds in five steps: initialize the
parameters with small random values; propagate an input forward to obtain a
prediction; compute the loss against the true output; backpropagate to obtain the
gradient of the loss with respect to every weight via the chain rule; and update
the weights with an optimization method. In practice the update is applied not to
the whole dataset at once but to small batches, giving stochastic gradient
descent.

\begin{definition}[Gradient descent]
\label{def:found:gd}
Let $\mathcal{L} : \R^p \to \R$ be a differentiable loss. \emph{Gradient
descent} is the iteration
\[
  \theta_{t+1} = \theta_t - \eta\, \nabla_\theta \mathcal{L}(\theta_t),
\]
where $\eta > 0$ is the \emph{learning rate}.
\end{definition}

\begin{definition}[Stochastic gradient descent]
\label{def:found:sgd}
Given a dataset $\{(x_i, y_i)\}_{i=1}^N$ and empirical loss
$\mathcal{L}(\theta) = \frac{1}{N}\sum_i \ell(f_\theta(x_i), y_i)$,
\emph{stochastic gradient descent} (SGD) approximates the gradient on a random
mini-batch $\mathcal{B} \subset \{1,\ldots,N\}$:
\[
  \theta_{t+1} = \theta_t - \eta \cdot \frac{1}{\abs{\mathcal{B}}}
  \sum_{i \in \mathcal{B}} \nabla_\theta \ell(f_\theta(x_i), y_i).
\]
\end{definition}

The mini-batch gradient is an \emph{unbiased} estimator of the true gradient,
$\E_{\mathcal{B}}[\,\cdot\,] = \nabla_\theta \mathcal{L}(\theta)$. Its higher
variance can act as implicit regularization, helping to escape suboptimal local
minima.

\subsection{Backpropagation}
\label{sec:found:backprop}

The backpropagation algorithm computes the gradients required by gradient descent
efficiently by exploiting the graph structure of the network.

\begin{definition}[Computational graph]
\label{def:found:compgraph}
A \emph{computational graph} is a directed acyclic graph $G = (V, E)$ in which
each node $v$ represents a variable or a differentiable operation $f_v$; the
input nodes (no incoming edges) correspond to parameters $\theta$ and data $x$;
each edge $(u, v)$ indicates that $v$ depends on $u$; and the single output node
produces the loss $\mathcal{L}$.
\end{definition}

The three graphs --- layer graph, neuron graph, computational graph --- describe
the same network at decreasing levels of abstraction; the computational graph
unfolds each neuron into its elementary operations, allowing the chain rule to be
applied systematically.

\begin{theorem}[Chain rule on computational graphs]
\label{thm:found:backprop}
Let $G = (V, E)$ be a computational graph with output node $\mathcal{L}$. For
each intermediate node $v$ with successors $\{u_1, \ldots, u_k\}$,
\[
  \frac{\partial \mathcal{L}}{\partial v}
  = \sum_{i=1}^k \frac{\partial \mathcal{L}}{\partial u_i}\,
  \frac{\partial u_i}{\partial v}.
\]
Backpropagation computes all gradients $\partial\mathcal{L}/\partial v$ in time
$O(\abs{V} + \abs{E})$, linear in the size of the graph.
\end{theorem}

\begin{proof}
Since $G$ is a DAG, \Cref{thm:found:topological} yields a topological ordering
$v_1, \ldots, v_n$, computable in $O(\abs{V} + \abs{E})$. In the \emph{forward
pass} we traverse the nodes in this order; each operation node $v_i$ with
predecessors $\{u_1, \ldots, u_k\}$ is evaluated as $f_{v_i}(u_1, \ldots, u_k)$,
all predecessors having smaller index and thus already computed, at total cost
$O(\abs{V} + \abs{E})$. For the \emph{backward pass}, the loss depends on an
intermediate node $v$ only through its successors, so the multivariable chain
rule gives the stated formula. We traverse the nodes in reverse topological
order, initializing $\partial\mathcal{L}/\partial\mathcal{L} = 1$ and computing,
for each $v_i$, $\partial\mathcal{L}/\partial v_i = \sum_{(v_i, u_j) \in E}
(\partial\mathcal{L}/\partial u_j)(\partial u_j/\partial v_i)$, using gradients
of successors already available. Each edge contributes once and each local
derivative is $O(1)$, so the backward pass is $O(\abs{V} + \abs{E})$, and so is
the whole algorithm.
\end{proof}

In practice modern frameworks implement \emph{automatic
differentiation}~\citep{Neidinger2010IntroductionTA}: during the forward pass
they build the computational graph dynamically, and during the backward pass they
traverse it in reverse topological order applying the chain rule, freeing the
user from computing derivatives by hand.

\subsection{The vanishing gradient problem}
\label{sec:found:vanishing}

A direct consequence of the chain rule is the \emph{vanishing gradient problem}.
By \Cref{thm:found:backprop}, the gradient of the loss with respect to the
parameters of the first layers is proportional to a product of activation
derivatives across the network,
\[
  \frac{\partial \mathcal{L}}{\partial \theta^{(1)}} \propto
  \prod_{\ell=1}^L \sigma'(z^{(\ell)}).
\]
If each $\abs{\sigma'(z^{(\ell)})}$ is below $1$ (as with saturating functions
like the sigmoid, where $\abs{\sigma'(z)} \le 0.25$), this product decreases
exponentially with depth $L$, so the gradients of the first layers become
numerically zero and those layers barely learn. This was a principal obstacle to
training deep networks in the 1990s. The solution came with non-saturating
activations such as ReLU, where $\sigma'(z) = 1$ for $z > 0$, letting gradients
flow without attenuation (\Cref{def:found:nonsaturating}).

\subsection{The universal approximation theorem}
\label{sec:found:universal}

The universal approximation theorem guarantees that any continuous function can
be approximated arbitrarily well by a network with a single hidden layer. Its
proof rests on tools from measure theory and functional analysis --- the notions
of measure and $\sigma$-algebra underpinning the representation theorem of Riesz
and the separation theorem of Hahn--Banach. To keep the focus on the result, we
collect those prerequisites in the appendix (\Cref{ch:appendix}) and state here
only what is needed.

\begin{definition}[Sigmoidal function]
\label{def:found:sigmoidal}
A function $\phi : \R \to \R$ is \emph{sigmoidal} if it is continuous, bounded,
and satisfies $\lim_{t\to-\infty}\phi(t) = 0$ and
$\lim_{t\to+\infty}\phi(t) = 1$. Classical examples are the logistic function and
a rescaled hyperbolic tangent.
\end{definition}

\begin{theorem}[Universal approximation; Cybenko, 1989]
\label{thm:found:universal}
Let $\phi$ be a sigmoidal function and $I_m = [0,1]^m$ the unit cube. The set of
functions
\[
  F(x) = \sum_{i=1}^N v_i\, \phi(w_i^\top x + b_i), \qquad
  N \in \mathbb{N},\ v_i, b_i \in \R,\ w_i \in \R^m,
\]
is dense in $C(I_m)$ in the uniform norm: for every $f \in C(I_m)$ and every
$\epsilon > 0$ there are parameters with
$\sup_{x \in I_m} \abs{F(x) - f(x)} < \epsilon$.
\end{theorem}

The restriction to $I_m$ is no real limitation: any compact $K \subset \R^m$ is
bounded (Heine--Borel, \Cref{ch:appendix}) and an affine map carries the problem
to $I_m$, so the result extends to $C(K)$. The theorem has been extended to any
non-polynomial activation~\citep{leshno1993multilayer}, including ReLU, where a
single hidden layer suffices with arbitrary width while bounded-width networks
require sufficient depth~\citep{lu2017expressive,hanin2019universal}; to $L^p$
norms~\citep{hornik1991approximation}; and to approximation-rate bounds on the
number of neurons~\citep{barron1993universal}.

\begin{proof}
We follow \citet{cybenko1989} via functional-analytic duality. The strategy is
indirect: rather than constructing the approximation, we show that no continuous
function can escape being approximated. Were the networks not dense in $C(I_m)$,
the Hahn--Banach corollary (\Cref{ch:appendix}) would furnish a nonzero
continuous linear functional vanishing on all of them; by the Riesz
representation theorem (\Cref{ch:appendix}) this is a nonzero finite signed
measure invisible to all neurons. The heart of the proof is to show no such
measure exists, i.e.\ that $\phi$ is \emph{discriminatory}.

\emph{Step 1: discriminatory functions.} Call $\phi$ \emph{discriminatory} if,
for every finite signed measure $\mu$ on $I_m$,
\[
  \int_{I_m} \phi(w^\top x + b)\, d\mu(x) = 0 \ \ \forall w, b
  \implies \mu = 0.
\]

\emph{Step 2: sigmoidal functions are discriminatory.} Suppose $\mu$ satisfies
the hypothesis. For $w, b$ and an auxiliary $\alpha$, set
$\psi_\lambda(x) = \phi(\lambda(w^\top x + b) + \alpha)$. As $\lambda \to \infty$,
$\psi_\lambda$ converges pointwise to $1$ on the open half-space
$H^+_{w,b} = \{w^\top x + b > 0\}$, to $0$ on $\{w^\top x + b < 0\}$, and to
$\phi(\alpha)$ on the hyperplane $H_{w,b} = \{w^\top x + b = 0\}$. Since
$\abs{\psi_\lambda} \le \sup_t \abs{\phi(t)} < \infty$ uniformly and this bound is
integrable over the compact $I_m$, the dominated convergence theorem
(\Cref{ch:appendix}) gives
\[
  \int_{I_m}\psi_\lambda\, d\mu \to \mu(H^+_{w,b}) + \phi(\alpha)\,\mu(H_{w,b}).
\]
The left side is $0$ for all $\lambda$, so
$\mu(H^+_{w,b}) + \phi(\alpha)\,\mu(H_{w,b}) = 0$ for all $\alpha$. This is affine
in $\phi(\alpha)$, and since $\phi$ takes a continuous range of values, both
$\mu(H_{w,b}) = 0$ and $\mu(H^+_{w,b}) = 0$. Thus $\mu$ vanishes on every open
half-space. As the half-spaces generate the Borel $\sigma$-algebra and a finite
signed measure is determined by its values on a generating class stable under
intersection, $\mu = 0$.

\emph{Step 3: density.} Let $\mathcal{S}$ be the set of functions
$G(x) = \sum_i v_i \phi(w_i^\top x + b_i)$. By the Hahn--Banach corollary
(\Cref{ch:appendix}), $\mathcal{S}$ is dense in $C(I_m)$ if and only if the only
continuous linear functional vanishing on $\mathcal{S}$ is zero. Let $L$ be such
a functional; by the Riesz representation theorem there is a finite signed
measure $\mu$ with $L(h) = \int_{I_m} h\, d\mu$. Then $L(G) = 0$ for all $G$
gives $\int_{I_m}\phi(w^\top x + b)\, d\mu = 0$ for all $w, b$, so by Step 2
$\mu = 0$ and $L = 0$. Hence $\mathcal{S}$ is dense, which is the assertion.
\end{proof}

This guarantees that a single-hidden-layer network can approximate any continuous
function arbitrarily well. It is, however, an \emph{existence} result, not a
\emph{constructive} one: it does not say how to find the weights, the required
number of neurons may be prohibitively large, and deep networks can represent
complex functions far more efficiently (\Cref{sec:found:depth}).

\subsection{Convergence theory}
\label{sec:found:convergence}

The universal approximation theorem guarantees that a weight configuration
exists; gradient descent is the mechanism that searches for it. We now state the
conditions under which this search converges. The proofs use a regularity
property, $L$-smoothness, whose equivalent characterizations and the convexity
inequality are recorded in \Cref{ch:appendix}.

\begin{definition}[$L$-smoothness and convexity]
\label{def:found:smooth}
A differentiable $f : \R^n \to \R$ is \emph{$L$-smooth} if
$\norm{\nabla f(x) - \nabla f(y)} \le L\norm{x - y}$ for all $x, y$. It is
\emph{convex} if $f(\lambda x + (1-\lambda)y) \le \lambda f(x) + (1-\lambda)f(y)$
for $\lambda \in [0,1]$; for differentiable $f$ this is equivalent to
$f(y) \ge f(x) + \inner{\nabla f(x)}{y - x}$ for all $x, y$.
\end{definition}

\begin{theorem}[Convergence of gradient descent, convex case]
\label{thm:found:gd-convex}
Let $f$ be convex and $L$-smooth with global minimum $f^* = f(\theta^*)$.
Gradient descent with learning rate $\eta = 1/L$ satisfies
\[
  f(\theta_T) - f^* \le \frac{L\norm{\theta_0 - \theta^*}^2}{2T},
\]
i.e.\ it converges at rate $O(1/T)$.
\end{theorem}

\begin{proof}
The descent lemma from $L$-smoothness (\Cref{ch:appendix}) gives
$f(y) \le f(x) + \inner{\nabla f(x)}{y - x} + \frac{L}{2}\norm{y - x}^2$. With
$x = \theta_t$ and $y = \theta_{t+1} = \theta_t - \frac{1}{L}\nabla f(\theta_t)$,
$f(\theta_{t+1}) \le f(\theta_t) - \frac{1}{2L}\norm{\nabla f(\theta_t)}^2$, so
the iterates decrease. Expanding the distance to the optimum,
\[
  \norm{\theta_{t+1} - \theta^*}^2 = \norm{\theta_t - \theta^*}^2
  - \tfrac{2}{L}\inner{\nabla f(\theta_t)}{\theta_t - \theta^*}
  + \tfrac{1}{L^2}\norm{\nabla f(\theta_t)}^2.
\]
Convexity gives $\inner{\nabla f(\theta_t)}{\theta_t - \theta^*} \ge f(\theta_t)
- f^*$, and the descent inequality gives
$\norm{\nabla f(\theta_t)}^2 \le 2L(f(\theta_t) - f(\theta_{t+1}))$. Substituting,
\[
  f(\theta_{t+1}) - f^* \le \tfrac{L}{2}\bigl(\norm{\theta_t - \theta^*}^2
  - \norm{\theta_{t+1} - \theta^*}^2\bigr).
\]
Summing over $t = 0, \ldots, T-1$ telescopes to
$\sum_t (f(\theta_{t+1}) - f^*) \le \frac{L}{2}\norm{\theta_0 - \theta^*}^2$.
Since $\{f(\theta_t)\}$ is decreasing, $T(f(\theta_T) - f^*) \le \sum_t
(f(\theta_{t+1}) - f^*)$, giving the bound on dividing by $T$.
\end{proof}

\begin{theorem}[Convergence of SGD]
\label{thm:found:sgd-convergence}
Let $f$ be convex and $L$-smooth, and let $g_t$ be an unbiased gradient estimator
with bounded variance, $\E[g_t] = \nabla f(\theta_t)$ and
$\E[\norm{g_t - \nabla f(\theta_t)}^2] \le \sigma^2$, and bounded gradients
$\norm{\nabla f(\theta_t)} \le G$. With $\eta_t = \eta_0/\sqrt{t}$,
\[
  \E[f(\bar\theta_T)] - f^* = O\!\left(\frac{(\sigma^2 + G^2)\log T}{\sqrt{T}}
  + \frac{\norm{\theta_0 - \theta^*}^2}{\eta_0\sqrt{T}}\right),
\]
where $\bar\theta_T = \frac{1}{T}\sum_{t=1}^T \theta_t$.
\end{theorem}

\begin{proof}[Sketch]
Following~\citet{bottou2018optimization}, write $g_t = \nabla f(\theta_t) +
\xi_t$ with $\E[\xi_t \mid \theta_t] = 0$ and
$\E[\norm{\xi_t}^2 \mid \theta_t] \le \sigma^2$. Taking conditional expectation
of $\norm{\theta_{t+1} - \theta^*}^2 = \norm{\theta_t - \theta^*}^2 -
2\eta_t \inner{g_t}{\theta_t - \theta^*} + \eta_t^2\norm{g_t}^2$ and using
convexity gives
$\E[\norm{\theta_{t+1} - \theta^*}^2 \mid \theta_t] \le \norm{\theta_t -
\theta^*}^2 - 2\eta_t(f(\theta_t) - f^*) + \eta_t^2(\norm{\nabla f(\theta_t)}^2 +
\sigma^2)$. Summing telescopes the distance terms; with $\eta_t = \eta_0/\sqrt{t}$
and $\sum_t 1/t \le 1 + \log T$, and bounding the left side below by
$\eta_T = \eta_0/\sqrt{T}$, we obtain
$\frac{1}{T}\sum_t \E[f(\theta_t) - f^*] \le (\norm{\theta_0 - \theta^*}^2 +
\eta_0^2(G^2 + \sigma^2)(1 + \log T))/(2\eta_0\sqrt{T})$. Jensen's inequality on
$\bar\theta_T$ gives the claim.
\end{proof}

The $O(\log T/\sqrt{T})$ rate of SGD is slower than the $O(1/T)$ of deterministic
gradient descent, but SGD processes far less data per iteration.

\begin{theorem}[Convergence to stationary points, non-convex case]
\label{thm:found:gd-nonconvex}
Let $f$ be $L$-smooth (not necessarily convex) and bounded below by $f^*$. With
$\eta \le 1/L$, gradient descent satisfies
\[
  \min_{t=0,\ldots,T-1} \norm{\nabla f(\theta_t)}^2 \le
  \frac{2(f(\theta_0) - f^*)}{\eta T},
\]
i.e.\ it finds an $\epsilon$-stationary point in $O(1/\epsilon^2)$ iterations.
\end{theorem}

\begin{proof}
The descent lemma gives $f(\theta_{t+1}) \le f(\theta_t) - \eta(1 -
\frac{L\eta}{2})\norm{\nabla f(\theta_t)}^2$; with $\eta \le 1/L$ the factor is
at least $\frac12$, so $f(\theta_{t+1}) \le f(\theta_t) -
\frac{\eta}{2}\norm{\nabla f(\theta_t)}^2$. Summing over $t = 0, \ldots, T-1$
telescopes to $\frac{\eta}{2}\sum_t \norm{\nabla f(\theta_t)}^2 \le f(\theta_0) -
f^* $ (using $f(\theta_T) \ge f^*$). Since the minimum of a finite collection
does not exceed its average, $\min_t \norm{\nabla f(\theta_t)}^2 \le
\frac{1}{T}\sum_t \norm{\nabla f(\theta_t)}^2 \le 2(f(\theta_0) - f^*)/(\eta T)$.
\end{proof}

Neural network losses are typically non-convex; \Cref{thm:found:gd-nonconvex}
guarantees convergence to stationary points, not global minima. In practice the
overparametrization of modern networks creates many equivalent minima, the noise
of SGD helps escape saddle points, and the minima found tend to generalize well.

\subsection{Advanced optimizers}
\label{sec:found:optimizers}

Modern optimizers improve on plain SGD by incorporating historical gradient
information~\citep{kingma2014adam,duchi2011adagrad}. Below $\odot$ denotes the
Hadamard product and $\oslash$ elementwise division. \emph{SGD with
momentum}~\citep{polyak1964methods} introduces a velocity
$v_t = \beta v_{t-1} + \nabla_\theta\mathcal{L}(\theta_{t-1})$,
$\theta_t = \theta_{t-1} - \eta v_t$, smoothing oscillations in high-curvature
directions. \emph{Adagrad}~\citep{duchi2011adagrad} adapts the learning rate per
parameter, $G_t = G_{t-1} + g_t \odot g_t$,
$\theta_t = \theta_{t-1} - \eta\, g_t \oslash (\sqrt{G_t} + \epsilon)$, so
historically large gradients receive smaller steps.
\emph{RMSProp}~\citep{hinton2012rmsprop} replaces the monotone accumulator by an
exponential moving average, $v_t = \beta v_{t-1} + (1-\beta)(g_t \odot g_t)$,
preventing the denominator from growing without bound. \emph{Adam}~\citep{kingma2014adam}
combines momentum with adaptive rates, maintaining estimates of the first and
second moments with bias correction:
\[
  m_t = \beta_1 m_{t-1} + (1-\beta_1)g_t, \quad
  v_t = \beta_2 v_{t-1} + (1-\beta_2)(g_t \odot g_t),
\]
$\hat m_t = m_t/(1-\beta_1^t)$, $\hat v_t = v_t/(1-\beta_2^t)$,
$\theta_t = \theta_{t-1} - \eta\, \hat m_t \oslash (\sqrt{\hat v_t} + \epsilon)$,
with typical values $\beta_1 = 0.9$, $\beta_2 = 0.999$, $\epsilon = 10^{-8}$.

\section{Deep learning}
\label{sec:found:deep}

Interest in neural models began when Rosenblatt showed the perceptron could learn
to classify patterns from data~\citep{rosenblatt1958perceptron}; Minsky and
Papert then proved the simple perceptron can only represent linearly separable
functions~\citep{minsky1969perceptrons}, prompting a long decline in interest.
Backpropagation gave an efficient way to train multilayer networks, but for
decades deep networks remained impractical due to the vanishing gradient problem
(\Cref{sec:found:vanishing}). \emph{Deep learning} emerged as the convergence of
several advances: initialization schemes that preserve activation variance across
layers, non-saturating activations such as ReLU, regularization methods such as
dropout, and specialized architectures that exploit the structure of the data ---
for example, convolutional networks exploit spatial locality to require orders of
magnitude fewer parameters than equivalent dense networks.

Deeper networks have greater expressive power but require more data. A common
remedy for scarce data is \emph{data augmentation}.

\begin{definition}[Data augmentation]
\label{def:found:augmentation}
Let $\{(\mathbf{x}_i, y_i)\}_{i=1}^N$ be a training set and $G$ a group of
transformations acting on the input space $\mathcal{X}$ preserving labels, so the
correct label of $g \cdot \mathbf{x}$ equals that of $\mathbf{x}$ for all
$g \in G$. \emph{Data augmentation} replaces the dataset by
$\{(g \cdot \mathbf{x}_i, y_i) : i, g \in G\}$.
\end{definition}

Augmentation acts as implicit regularization, forcing the model to learn
representations invariant to $G$. As discussed in \Cref{ch:priors}, this strategy
is inefficient compared with the equivariant architectures of geometric deep
learning, which build the symmetries directly into their structure.

\subsection{Advantages of depth}
\label{sec:found:depth}

Deep networks offer several fundamental advantages over wide shallow ones. The
clearest is \emph{representation efficiency}: there exist functions exhibiting an
exponential separation between deep and shallow networks. The parity function
$f(x_1, \ldots, x_n) = \bigoplus_i x_i$ requires $O(2^n)$ neurons in a two-layer
network but only $O(n)$ in one of $O(\log n)$ layers.

\begin{theorem}[Exponential separation~\citep{telgarsky2016benefitsdepthneuralnetworks}]
\label{thm:found:telgarsky}
For every integer $k \ge 1$ there exist functions computable by networks of depth
$\Theta(k^3)$ with $\Theta(1)$ neurons per layer that any network of depth
$O(k)$ would need $\Omega(2^k)$ neurons to approximate with constant error.
\end{theorem}

\begin{proof}[Sketch]
Define the triangle function with ReLU neurons,
\[
  \varphi(x) = 2\ReLU(x) - 4\ReLU(x - \tfrac12) + 2\ReLU(x - 1),
\]
a ``tent'' on $[0,1]$ peaking at $x = \tfrac12$. Let $\varphi_1 = \varphi$ and
$\varphi_{k+1} = \varphi \circ \varphi_k$. Each composition doubles the number of
oscillations: $\varphi_k$ has exactly $2^k$ oscillations on $[0,1]$, since
$\varphi$ maps $[0,1] \to [0,1]$ and each arc generates two. Since each
composition uses $O(1)$ extra neurons, $\varphi_k$ is computable by a network of
depth $O(k)$ and constant width. A ReLU network with $p$ neurons produces a
piecewise-linear function with at most $O(p^L)$ segments for depth $L$, so a
depth-$O(k)$ network needs $\Omega(2^k)$ neurons to match $2^k$ oscillations.
Refined segment-counting yields the $\Theta(k^3)$ versus $O(k)$ statement; see
\citet{telgarsky2016benefitsdepthneuralnetworks}.
\end{proof}

Depth further enables \emph{hierarchical learning} along the natural
compositional structure of many problems (pixels $\to$ edges $\to$ parts $\to$
objects); \emph{feature reuse}, sharing computation across processing paths so
that $m$ functions cost $O(n + m)$ parameters rather than $O(mn)$; and an
implicit inductive bias towards simpler functions, which tends to generalize
well.

\subsection{Modern architectures}
\label{sec:found:modern}

Deep learning has produced specialized architectures, each exploiting a specific
data structure. \emph{Feedforward} networks, built from dense layers, are the
basic architecture with unidirectional flow $f(x) = f_L \circ \cdots \circ
f_1(x)$. \emph{Convolutional} networks (\Cref{ch:grids}) exploit spatial locality
and translational invariance through convolution, drastically reducing parameters
via weight sharing and achieving translational equivariance. \emph{Recurrent}
networks process sequences through a hidden state, with LSTM and GRU variants
overcoming the vanishing gradient. \emph{Attention mechanisms} model long-range
dependencies without recurrence by dynamically weighting the contributions of
different parts of the input.

\begin{definition}[Self-attention layer]
\label{def:found:attention}
Let $\{h_1, \ldots, h_n\} \subset \R^d$ be a sequence of input vectors. A
\emph{self-attention} layer with key dimension $d_k$ uses three linear
projections $Q, K, V \in \R^{d \times d_k}$ and computes
\[
  \mathrm{SelfAttention}(h_i) = \sum_{j=1}^n \alpha_{ij}\, V h_j, \qquad
  \alpha_{ij} = \frac{\exp\bigl((Qh_i)^\top(Kh_j)/\sqrt{d_k}\bigr)}
  {\sum_{m=1}^n \exp\bigl((Qh_i)^\top(Kh_m)/\sqrt{d_k}\bigr)},
\]
where the \emph{attention weights} $\alpha_{ij} \in (0,1)$ satisfy
$\sum_j \alpha_{ij} = 1$.
\end{definition}

This formulation is invariant to permutations of the inputs and generalizes
pooling and convolution; it is the basis of the \emph{transformer} architecture
(\Cref{ch:graphs}). \emph{Graph neural networks} generalize neural operations to
non-Euclidean data through message passing, a central theme of \Cref{ch:graphs};
\emph{generative adversarial networks} and \emph{diffusion models} are generative
architectures. Many of these reappear in later chapters as instances of a single
geometric construction.

With the neural foundations in place, we now develop the geometric language that
geometric deep learning rests on: manifolds, groups and their representations,
and homogeneous spaces.

\section{Manifolds and tangent spaces}
\label{sec:found:manifolds}

A surface, a configuration space, or the set of rotations is not a vector space,
yet each looks like flat Euclidean space when viewed closely enough. The notion
of a manifold makes this precise, and the tangent space provides the
linearization on which calculus --- and ultimately learning --- proceeds. We
assume the standard topological prerequisites (recorded in \Cref{ch:appendix}): a
\emph{topological space} is a set with a collection of open sets closed under
arbitrary unions and finite intersections; it is \emph{Hausdorff} if distinct
points have disjoint neighbourhoods; it is \emph{second countable} if its
topology has a countable base; and a \emph{homeomorphism} is a continuous
bijection with continuous inverse.

\subsection{Topological and differentiable manifolds}
\label{sec:found:diff-manifolds}

\begin{definition}[Topological manifold]
\label{def:found:topmanifold}
A \emph{topological manifold} of dimension $m$ is a Hausdorff, second-countable
space $\mathcal{M}$ that is \emph{locally Euclidean}: every point has an open
neighbourhood $U$ homeomorphic, via a \emph{chart} $\phi : U \to V \subseteq
\R^m$, to an open subset of $\R^m$.
\end{definition}

Intuitively a manifold is a space that locally looks like $\R^m$ even when its
global shape is very different: the Earth's surface is curved yet looks flat to a
pedestrian; a crumpled sheet of paper is locally two-dimensional; a robot arm's
configuration space is non-Euclidean but locally coordinatized by joint angles.

\begin{definition}[Atlas and charts]
\label{def:found:atlas}
An \emph{atlas} of dimension $m$ for $\mathcal{M}$ is a collection
$\mathcal{A} = \{(U_\alpha, \phi_\alpha)\}_{\alpha \in I}$ where
$\{U_\alpha\}$ is an open cover and each $\phi_\alpha : U_\alpha \to V_\alpha
\subseteq \R^m$ is a chart. The \emph{transition maps} $\phi_\beta \circ
\phi_\alpha^{-1} : \phi_\alpha(U_\alpha \cap U_\beta) \to \phi_\beta(U_\alpha \cap
U_\beta)$ describe how overlapping charts relate.
\end{definition}

\begin{figure}[h]
\centering
\begin{tikzpicture}[>=Stealth, scale=1.0]
\shade[ball color=gdlblue!18, opacity=0.85]
  (-2.4,-0.9) .. controls (-1.6,1.2) and (1.0,1.6) .. (2.4,0.6)
  .. controls (1.8,-1.2) and (-1.2,-1.7) .. (-2.4,-0.9);
\draw[thick, gdlblue!55]
  (-2.4,-0.9) .. controls (-1.6,1.2) and (1.0,1.6) .. (2.4,0.6)
  .. controls (1.8,-1.2) and (-1.2,-1.7) .. (-2.4,-0.9);
\node[text=gdlblue!60!black] at (-2.1,1.0) {$\mathcal{M}$};
\draw[gdlgreen!60!black, thick] (-1.5,0.2) ellipse (0.8cm and 0.55cm);
\node[font=\tiny, gdlgreen!50!black] at (-1.5,0.9) {$U_\alpha$};
\draw[gdlorange!70, thick] (0.9,0.3) ellipse (0.8cm and 0.55cm);
\node[font=\tiny, gdlorange!60!black] at (0.9,1.0) {$U_\beta$};
\draw[fill=gdlgreen!10, draw=gdlgreen!55] (-4.6,-2.6) rectangle (-2.9,-1.2);
\node[font=\tiny] at (-3.75,-1.0) {$V_\alpha \subseteq \R^m$};
\draw[fill=gdlorange!10, draw=gdlorange!60] (2.9,-2.6) rectangle (4.6,-1.2);
\node[font=\tiny] at (3.75,-1.0) {$V_\beta \subseteq \R^m$};
\draw[->, gdlgreen!55!black, thick] (-1.7,-0.2) to[bend right=15] node[left, font=\tiny] {$\phi_\alpha$} (-3.6,-1.3);
\draw[->, gdlorange!65, thick] (1.1,-0.1) to[bend left=15] node[right, font=\tiny] {$\phi_\beta$} (3.6,-1.3);
\draw[<->, gdlpurple!70, thick] (-2.85,-1.9) -- node[above, font=\tiny] {$\phi_\beta\circ\phi_\alpha^{-1}$} (2.85,-1.9);
\end{tikzpicture}
\caption[Atlas and charts of a manifold]{An atlas
$\mathcal{A} = \{(U_\alpha, \phi_\alpha)\}$ of a manifold $\mathcal{M}$: each
chart is a homeomorphism $\phi_\alpha : U_\alpha \to V_\alpha \subseteq \R^m$
defining local coordinates, and the transition map connects overlapping
charts.}
\label{fig:found:atlas}
\end{figure}

\begin{definition}[Differentiable manifold]
\label{def:found:diffmanifold}
A function $f : U \subseteq \R^n \to \R^m$ is \emph{smooth} (of class
$C^\infty$) if it has continuous partial derivatives of all orders. A topological
manifold $\mathcal{M}$ is \emph{differentiable} (smooth) if it has an atlas all
of whose transition maps are $C^\infty$. This is exactly the extra structure
needed to do calculus on $\mathcal{M}$.
\end{definition}

\subsection{The tangent space}
\label{sec:found:tangent}

The central object is the tangent space, which captures the infinitesimal
directions at a point. Two complementary constructions make it precise.

\begin{definition}[Tangent space via curves]
\label{def:found:tangent-curves}
Let $\mathcal{M}$ be smooth and $p \in \mathcal{M}$. Two curves
$\gamma_1, \gamma_2 : (-\epsilon,\epsilon) \to \mathcal{M}$ with
$\gamma_1(0) = \gamma_2(0) = p$ are \emph{tangent at $p$} if, for some (hence
every) chart $(U,\phi)$ around $p$, $(\phi\circ\gamma_1)'(0) =
(\phi\circ\gamma_2)'(0)$. The \emph{tangent space} $T_p\mathcal{M}$ is the set of
equivalence classes $[\gamma]$ of curves tangent at $p$.
\end{definition}

\begin{remark}[Vector-space structure]
\label{rem:found:tangent-vector}
The map $[\gamma] \mapsto (\phi\circ\gamma)'(0) \in \R^m$ is a bijection of
$T_p\mathcal{M}$ with $\R^m$, transporting the vector-space structure of $\R^m$
to $T_p\mathcal{M}$. This is independent of the chart: the chain rule applied to
the transition map guarantees that the isomorphisms induced by two charts are
compatible.
\end{remark}

\begin{definition}[Tangent space via derivations]
\label{def:found:tangent-deriv}
A \emph{derivation} at $p$ is a linear map $v : C^\infty(\mathcal{M}) \to \R$
obeying the Leibniz rule $v(fg) = v(f)\,g(p) + f(p)\,v(g)$. The tangent space
$T_p\mathcal{M}$ is the vector space of all derivations at $p$.
\end{definition}

The first construction reads a tangent vector as the initial velocity of a curve
--- a direction of motion; the second reads it as a directional-derivative
operator --- a direction of change of functions.

\begin{lemma}[Coordinate decomposition of derivations]
\label{lem:found:derivations}
Let $\mathcal{M}$ be a smooth $m$-manifold, $p \in \mathcal{M}$, and $(U, \phi)$ a
chart with $\phi(p) = 0$ and coordinate functions $x^1, \ldots, x^m$. Then
derivations are local (if $f = g$ near $p$ then $v(f) = v(g)$), and every
derivation $v$ at $p$ decomposes as
\[
  v = \sum_{j=1}^m v(x^j)\,\frac{\partial}{\partial x^j}\Big|_p,
  \qquad \frac{\partial}{\partial x^j}\Big|_p(f) =
  \frac{\partial(f\circ\phi^{-1})}{\partial x^j}\Big|_{\phi(p)},
\]
so $\dim T_p\mathcal{M} = m$.
\end{lemma}

\begin{proof}
\emph{Locality.} Let $\chi$ be a smooth bump function with $\chi \equiv 1$ near
$p$ and support in $U$. If $f = g$ on $U$ then $\chi(f-g) = 0$ on $\mathcal{M}$,
so by Leibniz $0 = v(\chi(f-g)) = v(\chi)(f-g)(p) + \chi(p)v(f-g) = v(f-g)$, since
$\chi(p) = 1$ and $(f-g)(p) = 0$.

\emph{Decomposition.} First $v$ annihilates constants: $v(1) = v(1\cdot 1) =
2v(1)$ gives $v(1) = 0$, hence $v(c) = 0$. Writing $\hat f = f\circ\phi^{-1}$, the
fundamental theorem of calculus with integral remainder gives, on a convex
neighbourhood of $0$,
$\hat f(\mathbf{x}) = \hat f(0) + \sum_j x^j g_j(\mathbf{x})$ with $g_j$ smooth
and $g_j(0) = \partial\hat f/\partial x^j(0)$. Composing with $\phi$ and using
locality, $f = f(p) + \sum_j x^j (g_j\circ\phi)$ on $U$. Applying $v$, using that
it annihilates constants, $x^j(p) = 0$, and Leibniz, $v(f) = \sum_j v(x^j)
\partial(f\circ\phi^{-1})/\partial x^j|_0$. Since
$\frac{\partial}{\partial x^j}|_p(x^k) = \delta_j^k$, the operators
$\{\partial/\partial x^j|_p\}$ are independent and span, so the dimension is $m$
\citep[Prop.~3.2]{lee2012introduction}.
\end{proof}

\begin{proposition}[Equivalence of the two constructions]
\label{prop:found:tangent-equiv}
The two definitions of tangent space are canonically isomorphic via
$\Phi([\gamma])(f) = \frac{d}{dt}\big|_0 (f\circ\gamma)(t)$, the isomorphism
being independent of coordinates.
\end{proposition}

\begin{proof}
\emph{Well defined.} If $\gamma_1, \gamma_2$ are tangent at $p$, a chart
$(U,\phi)$ with coordinates $x^j$ gives, by the chain rule,
$(f\circ\gamma_i)'(0) = \sum_j \partial(f\circ\phi^{-1})/\partial x^j|_{\phi(p)}
\cdot (x^j\circ\gamma_i)'(0)$; since the velocities agree, so do the derivatives.
\emph{Derivation.} Linearity is immediate, and the product rule for
$t \mapsto f(\gamma(t))$ gives Leibniz, $v_\gamma(fg) = v_\gamma(f)g(p) +
f(p)v_\gamma(g)$. \emph{Injective.} If $\Phi([\gamma_1]) = \Phi([\gamma_2])$, then
taking $f = x^j$ gives $(x^j\circ\gamma_1)'(0) = (x^j\circ\gamma_2)'(0)$ for all
$j$, so $(\phi\circ\gamma_1)'(0) = (\phi\circ\gamma_2)'(0)$ and the curves are
tangent. \emph{Surjective.} Given a derivation $v$, set $a^j = v(x^j)$ and
$\gamma(t) = \phi^{-1}(t\mathbf{a})$; then $\Phi([\gamma]) = \sum_j v(x^j)
\partial/\partial x^j|_p = v$ by \Cref{lem:found:derivations}.
\end{proof}

\begin{example}[Tangent space to the sphere]
\label{ex:found:tangent-sphere}
At the north pole $p = (0,0,1)$ of $S^2$, the curve
$\gamma(t) = (\sin t, 0, \cos t)$ has $\gamma'(0) = (1,0,0)$. The tangent space
$T_p S^2$ is the horizontal plane $\{(x,y,0)\}$, reflecting $\dim S^2 = 2$.
\end{example}

\begin{figure}[h]
\centering
\begin{tikzpicture}[>=Stealth, scale=1.0]
\shade[ball color=gdlblue!18, opacity=0.9]
    (-2,-1.3) .. controls (-1,1) and (1,2) .. (2.2,0.6)
    .. controls (1.6,-1) and (-1,-1.9) .. (-2,-1.3);
\draw[thick, gdlblue!55]
    (-2,-1.3) .. controls (-1,1) and (1,2) .. (2.2,0.6)
    .. controls (1.6,-1) and (-1,-1.9) .. (-2,-1.3);
\node[text=gdlblue!60!black, font=\small] at (-1.7,-1.5) {$\mathcal{M}$};
\coordinate (p) at (0,0.35);
\begin{scope}[on background layer]
\fill[gdlorange!18, opacity=0.85]
    ($(p)+(-1.7,-0.55)$) -- ($(p)+(1.7,0.35)$) --
    ($(p)+(1.4,1.15)$) -- ($(p)+(-2.0,0.25)$) -- cycle;
\draw[gdlorange!60]
    ($(p)+(-1.7,-0.55)$) -- ($(p)+(1.7,0.35)$) --
    ($(p)+(1.4,1.15)$) -- ($(p)+(-2.0,0.25)$) -- cycle;
\end{scope}
\node[text=gdlorange!60!black, font=\small] at (1.9,1.25) {$T_p\mathcal{M}$};
\draw[thick, gdlpurple!75] (-1.4,-0.4) .. controls (-0.5,0.05) .. (p)
    .. controls (0.5,0.65) .. (1.4,1.0);
\node[gdlpurple!75, font=\tiny] at (1.55,1.05) {$\gamma_1$};
\draw[thick, gdlgreen!60!black] (-0.5,-0.9) .. controls (-0.2,-0.25) .. (p)
    .. controls (0.2,0.9) .. (0.45,1.4);
\node[gdlgreen!60!black, font=\tiny] at (0.6,1.45) {$\gamma_2$};
\draw[->, very thick, gdlpurple!75] (p) -- ++(0.95,0.45)
    node[right, font=\tiny] {$\gamma_1'(0)$};
\draw[->, very thick, gdlgreen!60!black] (p) -- ++(0.25,0.85)
    node[above, font=\tiny] {$\gamma_2'(0)$};
\fill[gdlred!80] (p) circle (2.2pt);
\node[below right=0pt and 1pt of p, font=\small] {$p$};
\end{tikzpicture}
\caption[Tangent space]{The tangent space $T_p\mathcal{M}$ at a point $p$.
Tangent vectors are the initial velocities $\gamma'(0)$ of curves through $p$;
equivalently, directional-derivative operators on smooth functions.}
\label{fig:found:tangent}
\end{figure}

A smooth map $F : \mathcal{M} \to \mathcal{N}$ acts on tangent vectors through its
\emph{differential} (or pushforward) $dF_p : T_p\mathcal{M} \to
T_{F(p)}\mathcal{N}$, defined on derivations by $dF_p(v)(g) = v(g\circ F)$; in
local coordinates $dF_p$ is the Jacobian matrix. The map $F$ is an
\emph{immersion} if $dF_p$ is injective everywhere
($\rank dF_p = \dim\mathcal{M}$), a \emph{submersion} if it is surjective
everywhere ($\rank dF_p = \dim\mathcal{N}$), and an \emph{embedding} if it is an
injective immersion that is a homeomorphism onto its image. Immersion and
submersion are dual notions; a submersion has each fibre $F^{-1}(q)$ a smooth
submanifold of dimension $\dim\mathcal{M} - \dim\mathcal{N}$, a fact used for the
matrix Lie groups of \Cref{sec:found:lie-groups}. If $\mathcal{M}$ is compact,
every injective immersion into a Hausdorff space is automatically an embedding.

\subsection{Cotangent space and differential forms}
\label{sec:found:cotangent}

The tangent space captures \emph{directions} at a point; its dual, the
\emph{cotangent space}, captures how \emph{functions change} along those
directions. This passage from vectors to covectors underlies the Riemannian
metric, connections on bundles, and curvature developed in
\Cref{ch:geodesics,ch:gauges}.

\begin{definition}[Cotangent space and the differential]
\label{def:found:cotangent}
The \emph{cotangent space} $T_p^*\mathcal{M} = (T_p\mathcal{M})^*$ is the dual of
the tangent space; its elements are \emph{covectors}, and $\dim T_p^*\mathcal{M}
= m$. For $f \in C^\infty(\mathcal{M})$, the \emph{differential} $(df)_p \in
T_p^*\mathcal{M}$ is $(df)_p(v) = v(f)$. In local coordinates $df = \sum_i
\frac{\partial f}{\partial x^i}\, dx^i$, and the covectors $\{dx^1, \ldots,
dx^m\}$ form the \emph{dual basis}, characterized by $dx^i(\partial/\partial
x^j) = \delta^i_j$.
\end{definition}

The notation $dx^i$ denotes not an infinitesimal but a well-defined linear
functional: the covector extracting the $i$-th component of a tangent vector. The
\emph{tensor product} of covectors $\alpha, \beta \in T_p^*\mathcal{M}$ is the
bilinear form $(\alpha \otimes \beta)(v, w) = \alpha(v)\beta(w)$; the Riemannian
metric is written $g_p = \sum_{i,j} g_{ij}(p)\, dx^i \otimes dx^j$.

\begin{definition}[Differential forms and the exterior product]
\label{def:found:forms}
A \emph{$1$-form} is a smooth assignment $\omega : p \mapsto \omega_p \in
T_p^*\mathcal{M}$, written locally $\omega = \sum_i \omega_i(x)\, dx^i$; the space
of $1$-forms is $\Omega^1(\mathcal{M})$. A \emph{$k$-form} assigns to each point
an antisymmetric $(0,k)$-tensor; $\Omega^0(\mathcal{M}) = C^\infty(\mathcal{M})$
and $\Omega^k = 0$ for $k > m$. The \emph{exterior product} $\wedge :
\Omega^k \times \Omega^l \to \Omega^{k+l}$ is the normalized antisymmetrization
of the tensor product; for $1$-forms,
\[
  (\alpha \wedge \beta)(v, w) = \alpha(v)\beta(w) - \alpha(w)\beta(v),
\]
so $dx^i \wedge dx^j = -dx^j \wedge dx^i$ and $dx^i \wedge dx^i = 0$. In general
$\alpha \wedge \beta = (-1)^{kl}\beta \wedge \alpha$ (graded
anticommutativity), and every $k$-form is $\omega = \sum_{i_1 < \cdots < i_k}
\omega_{i_1\cdots i_k}(x)\, dx^{i_1} \wedge \cdots \wedge dx^{i_k}$.
\end{definition}

\begin{definition}[Exterior derivative]
\label{def:found:exterior-deriv}
The \emph{exterior derivative} is the unique linear operator
$d : \Omega^k(\mathcal{M}) \to \Omega^{k+1}(\mathcal{M})$ that extends the
differential on functions, satisfies the graded Leibniz rule $d(\alpha \wedge
\beta) = d\alpha \wedge \beta + (-1)^k \alpha \wedge d\beta$ for $\alpha \in
\Omega^k$, and is nilpotent, $d \circ d = 0$. Nilpotence yields the de Rham
complex $0 \to \Omega^0 \xrightarrow{d} \Omega^1 \xrightarrow{d} \cdots
\xrightarrow{d} \Omega^m \to 0$.
\end{definition}

\begin{example}[Forms on the sphere]
\label{ex:found:forms-sphere}
On $S^2$ with spherical coordinates $(\theta, \phi)$, the function
$f = \cos\theta$ is a $0$-form, its differential $df = -\sin\theta\, d\theta$ is a
$1$-form, and the area element $\sin\theta\, d\theta \wedge d\phi$ is the unique
$2$-form up to smooth scalar multiples, since $\dim S^2 = 2$. The area form
appears in spherical convolution, where signals on $S^2$ are integrated against
it.
\end{example}

These constructions reappear concretely in the rest of the book: the Riemannian
metric is a smooth field of tensor products of covectors; connections on
principal bundles are Lie-algebra-valued $1$-forms; the volume form
$\sqrt{\det(g_{ij})}\, dx^1 \wedge \cdots \wedge dx^m$ is the $m$-form that lets
us integrate signals over manifolds; and the de Rham complex is the continuous
counterpart of simplicial chains, connecting differential geometry with algebraic
topology (\Cref{ch:geodesics,ch:gauges}).

\section{Groups and group actions}
\label{sec:found:groups}

A symmetry is a transformation that leaves a quantity of interest unchanged. The
collection of all such transformations is never arbitrary: the composition of two
symmetries is a symmetry, every symmetry can be undone, and doing nothing is a
symmetry. These three facts are exactly the axioms of a group.

\begin{definition}[Group]
\label{def:found:group}
A \emph{group} is a set $\group$ with a binary operation
$\cdot : \group \times \group \to \group$ satisfying: (1) \emph{associativity},
$(a \cdot b) \cdot c = a \cdot (b \cdot c)$; (2) \emph{identity}, there is
$e \in \group$ with $e \cdot a = a \cdot e = a$; and (3) \emph{inverses}, each
$a$ has $a^{-1}$ with $a \cdot a^{-1} = a^{-1} \cdot a = e$. The group is
\emph{abelian} if $a \cdot b = b \cdot a$ for all $a, b$.
\end{definition}

Translations of an image form an abelian group; rotations in three dimensions do
not commute and their group is non-abelian. A subset $H \subseteq \group$ that is
a group under the inherited operation is a \emph{subgroup}, written
$H \le \group$; the subgroup criterion is that $H$ contains $e$, is closed under
the operation, and contains inverses. A subgroup $H$ is \emph{normal}
($H \triangleleft \group$) if $gHg^{-1} = H$ for every $g \in \group$.

The bridge between two groups is a map preserving their operations.

\begin{definition}[Homomorphism]
\label{def:found:homomorphism}
A map $\phi : \group \to \mathcal{H}$ is a \emph{homomorphism} if
$\phi(g_1 g_2) = \phi(g_1)\phi(g_2)$ for all $g_1, g_2$. A bijective
homomorphism is an \emph{isomorphism}, written $\group \cong \mathcal{H}$. The
\emph{kernel} $\ker(\phi) = \{g : \phi(g) = e_{\mathcal{H}}\}$ is always a normal
subgroup, and the \emph{image} $\Img(\phi) = \{\phi(g) : g \in \group\}$ is a
subgroup of $\mathcal{H}$.
\end{definition}

Isomorphic groups are algebraically indistinguishable, differing only in the
names of their elements. Homomorphisms are the formal device through which
abstract symmetries become computable transformations, and in
\Cref{ch:priors,ch:groups} they are what makes equivariance implementable.

To apply a group to images, graphs, or point clouds we must say how its elements
transform the points of a set.

\begin{definition}[Group action]
\label{def:found:action}
An \emph{action} of a group $\group$ on a set $X$ is a map
$\group \times X \to X$, $(g, x) \mapsto g \cdot x$, satisfying
$(g_1 g_2) \cdot x = g_1 \cdot (g_2 \cdot x)$ and $e \cdot x = x$ for all
$g_1, g_2 \in \group$ and $x \in X$.
\end{definition}

\begin{proposition}[Actions as homomorphisms]
\label{prop:found:action-hom}
Every action induces a homomorphism $\tilde\rho : \group \to \Sym(X)$,
$\tilde\rho(g)(x) = g \cdot x$, into the group of bijections of $X$.
\end{proposition}

\begin{proof}
Each $\tilde\rho(g)$ is a bijection with inverse $\tilde\rho(g^{-1})$, since
$g^{-1}\cdot(g\cdot x) = (g^{-1}g)\cdot x = x$. And $\tilde\rho(g_1 g_2)(x) =
(g_1 g_2)\cdot x = g_1\cdot(g_2\cdot x) = (\tilde\rho(g_1)\circ\tilde\rho(g_2))(x)$,
so $\tilde\rho$ preserves the operation.
\end{proof}

Studying how $\group$ acts on $X$ is therefore the same as studying
homomorphisms $\group \to \Sym(X)$. The actions relevant to this book are
familiar.

\begin{example}[Actions in deep learning]
\label{ex:found:actions}
\leavevmode
\begin{itemize}
  \item Translations $\Z^2$ act on discrete images by shifting pixels.
  \item Rotations $\SO(3)$ act on point clouds by matrix multiplication.
  \item Permutations $S_n$ act on graphs with $n$ nodes by relabelling.
\end{itemize}
These are the symmetry groups of grids, geometric data, and graphs, organizing
\Cref{ch:grids,ch:groups,ch:graphs}.
\end{example}

\subsection{Quotient groups and the first isomorphism theorem}
\label{sec:found:quotients}

When $H \triangleleft \group$ is normal, the left cosets $gH = \{gh : h \in H\}$
partition $\group$ and inherit a group structure.

\begin{definition}[Quotient group]
\label{def:found:quotient}
Let $H \triangleleft \group$. The \emph{quotient group} $\group/H$ is the set of
left cosets $\{gH : g \in \group\}$ with operation
$(g_1 H)(g_2 H) = (g_1 g_2) H$.
\end{definition}

\begin{proposition}[Well-definedness]
\label{prop:found:quotient-welldef}
The operation $(g_1 H)(g_2 H) = (g_1 g_2)H$ on $\group/H$ is well defined if and
only if $H \triangleleft \group$.
\end{proposition}

\begin{proof}
($\Leftarrow$) Suppose $H \triangleleft \group$ and $g_1' = g_1 h_1$,
$g_2' = g_2 h_2$ with $h_i \in H$. Then $g_1'g_2' = g_1 g_2 (g_2^{-1}h_1 g_2)h_2$,
and $g_2^{-1}h_1 g_2 \in H$ by normality, so $(g_1'g_2')H = (g_1 g_2)H$.
($\Rightarrow$) If the operation is well defined, then for $g \in \group$,
$h \in H$, $(ghg^{-1})H = (ghH)(g^{-1}H) = (gH)(g^{-1}H) = H$, so $ghg^{-1} \in H$,
i.e.\ $H \triangleleft \group$.
\end{proof}

\begin{theorem}[First isomorphism theorem]
\label{thm:found:first-iso}
For any homomorphism $\phi : \group \to \mathcal{H}$,
\[
  \group / \ker(\phi) \;\cong\; \Img(\phi), \qquad
  g\,\ker(\phi) \mapsto \phi(g).
\]
\end{theorem}

\begin{proof}
The map $\bar\phi(g\ker\phi) = \phi(g)$ is well defined: if $g_1\ker\phi =
g_2\ker\phi$ then $g_1^{-1}g_2 \in \ker\phi$, so $\phi(g_1) = \phi(g_2)$. It is a
homomorphism since $\bar\phi((g_1\ker\phi)(g_2\ker\phi)) = \phi(g_1 g_2) =
\phi(g_1)\phi(g_2)$. It is injective: $\bar\phi(g\ker\phi) = e$ gives
$g \in \ker\phi$, the identity coset. It is surjective onto $\Img(\phi)$ by
definition.
\end{proof}

Quotients model the systematic reduction of symmetry in hierarchical
architectures: a network on three-dimensional data may begin with full $\SE(3)$
symmetry and reduce to $\SO(3)$ in deeper layers through a homomorphism whose
kernel (the translations) is normal.

Many symmetry groups are built from simpler ones. The Euclidean group combines
translations with rotations and reflections, but not as a plain product: a
rotation reorients the translation that follows it.

\begin{definition}[Semidirect product]
\label{def:found:semidirect}
Given groups $N, H$ and a homomorphism $\phi : H \to \Aut(N)$, the
\emph{semidirect product} $N \rtimes_\phi H$ is the set $N \times H$ with product
\[
  (n_1, h_1)\cdot(n_2, h_2) = \bigl(n_1\,\phi(h_1)(n_2),\; h_1 h_2\bigr).
\]
\end{definition}

\begin{example}[Euclidean group]
\label{ex:found:euclidean}
The Euclidean group $E(n) = \R^n \rtimes \Orth(n)$ combines translations with
$\Orth(n)$ via $\phi(R)(t) = Rt$, so
$(t_1, R_1)\cdot(t_2, R_2) = (t_1 + R_1 t_2,\, R_1 R_2)$: first rotate, then
translate. Its orientation-preserving subgroup $\SE(n) = \R^n \rtimes \SO(n)$ is
the symmetry group of rigid three-dimensional data.
\end{example}

\subsection{Lie groups and Lie algebras}
\label{sec:found:lie-groups}

Permutations form a \emph{discrete} group; rotations and translations form
\emph{continuous} families. To do calculus over the continuous case --- to
differentiate and optimize over transformations --- we require the group to be a
smooth manifold.

\begin{definition}[Lie group]
\label{def:found:liegroup}
A \emph{Lie group} is a group $G$ that is also a smooth manifold and for which
multiplication $(g,h) \mapsto gh$ and inversion $g \mapsto g^{-1}$ are smooth. A
\emph{smooth action} of $G$ on a manifold $M$ is an action $G \times M \to M$
that is itself smooth.
\end{definition}

A Lie group is curved but linear to first order. Its tangent space at the
identity, $\Lie{g} = T_e G$, carries an antisymmetric bilinear bracket making it
a Lie algebra.

\begin{definition}[Lie algebra]
\label{def:found:liealgebra}
A \emph{Lie bracket} on a vector space $\Lie{g}$ is a bilinear, antisymmetric map
$[\cdot,\cdot] : \Lie{g} \times \Lie{g} \to \Lie{g}$ satisfying the Jacobi
identity $[X,[Y,Z]] + [Y,[Z,X]] + [Z,[X,Y]] = 0$. A vector space with a Lie
bracket is a \emph{Lie algebra}. Every Lie group $G$ carries the canonical Lie
algebra $\Lie{g} = T_e G$ with the bracket of left-invariant vector fields.
\end{definition}

The passage back from the algebra to the group is the exponential map.

\begin{definition}[Exponential map]
\label{def:found:exp}
The \emph{exponential map} $\exp : \Lie{g} \to G$ sends $X$ to $\gamma(1)$, where
$\gamma$ is the unique one-parameter subgroup with $\gamma'(0) = X$. For matrix
groups it is the matrix exponential $\exp(A) = \sum_{k\ge 0} A^k/k!$.
\end{definition}

\begin{example}[Exponential map for $\SO(2)$]
\label{ex:found:exp-so2}
The Lie algebra $\Lie{so}(2)$ is spanned by
$J = \left(\begin{smallmatrix}0 & -1\\ 1 & 0\end{smallmatrix}\right)$, and
\[
  \exp(\theta J) = \begin{pmatrix}\cos\theta & -\sin\theta\\
  \sin\theta & \cos\theta\end{pmatrix} = R_\theta,
\]
so the exponential turns an angle (an element of the algebra) into a rotation (an
element of the group).
\end{example}

\begin{figure}[h]
\centering
\begin{tikzpicture}[>=Stealth, scale=1.0]
\shade[ball color=gdlblue!14, opacity=0.55] (0,0) circle (1.7cm);
\draw[gdlblue!55, thick] (0,0) circle (1.7cm);
\node[gdlblue!60!black, font=\small] at (1.5,1.5) {$G$};
\fill[gdlred!80] (0,-1.7) circle (1.6pt) node[below, font=\tiny] {$e$};
\draw[fill=gdlorange!12, draw=gdlorange!60]
  (-2.3,-1.7) -- (2.3,-1.7) -- (2.9,-0.7) -- (-1.7,-0.7) -- cycle;
\node[gdlorange!60!black, font=\small] at (2.4,-1.05) {$\Lie{g} = T_e G$};
\draw[->, very thick, gdlpurple!75] (0,-1.7) -- (1.3,-1.25)
  node[midway, below, font=\tiny] {$X$};
\draw[->, thick, gdlgreen!60!black]
  (1.3,-1.25) .. controls (1.7,-0.6) and (1.7,0.2) .. (1.5,0.8)
  node[right, font=\tiny] {$\exp(X)$};
\fill[gdlgreen!60!black] (1.5,0.8) circle (1.4pt);
\end{tikzpicture}
\caption[Exponential map]{The exponential map $\exp : \Lie{g} \to G$ carries
vectors of the Lie algebra (the tangent space at the identity $e$) to elements of
the Lie group, connecting the linear structure of the algebra with the nonlinear
structure of the group.}
\label{fig:found:exp}
\end{figure}

This is why Lie algebras matter for network design: optimization can be carried
out in the flat vector space $\Lie{g}$ with ordinary gradients, and the resulting
transformation applied in the group through $\exp$.

\subsection{Fundamental Lie groups}
\label{sec:found:fundamental-groups}

The dimension of a Lie group is its dimension as a manifold, equal to the
dimension of its Lie algebra, $\dim G = \dim\Lie{g} = \dim T_e G$. Dimensions of
matrix groups follow from the regular level set theorem (\Cref{ch:appendix}): if
$F : \R^{n\times n} \to \R^k$ is smooth and $c$ is a regular value, then
$F^{-1}(c)$ is a submanifold of dimension $n^2 - k$.

\begin{definition}[Orthogonal and special orthogonal groups]
\label{def:found:orthogonal}
The \emph{orthogonal group} is $\Orth(n) = \{R \in \R^{n\times n} : R^\top R =
I_n\}$, containing rotations ($\det R = +1$) and reflections ($\det R = -1$). The
\emph{special orthogonal group} $\SO(n) = \{R \in \Orth(n) : \det R = 1\}$ of
rotations is a compact Lie group of dimension $\frac{n(n-1)}{2}$.
\end{definition}

The dimension follows by applying the regular level set theorem to
$F(R) = R^\top R - I_n$ valued in symmetric matrices: $dF_R(H) = H^\top R +
R^\top H$ is surjective (take $H = \frac12 RS$ to hit any symmetric $S$), so
$\dim\Orth(n) = n^2 - \frac{n(n+1)}{2} = \frac{n(n-1)}{2}$; the determinant
condition is discrete and selects the identity component. As a check, the Lie
algebra $\Lie{so}(n)$ of antisymmetric matrices has exactly $\binom{n}{2}$ free
entries.

\begin{example}[$\SO(2)$ and $\SO(3)$]
\label{ex:found:so2-so3}
A planar rotation is a single angle,
$R_\theta = \left(\begin{smallmatrix}\cos\theta & -\sin\theta\\ \sin\theta &
\cos\theta\end{smallmatrix}\right)$, so $\SO(2) \cong S^1$ is the circle, with
$\Lie{so}(2)$ generated by $J$. The group $\SO(3)$ has dimension three; its Lie
algebra $\Lie{so}(3)$ of antisymmetric $3\times 3$ matrices has basis
$L_x, L_y, L_z$ with $[L_i, L_j] = \epsilon_{ijk} L_k$. The exponential has the
closed form known as the Rodrigues formula: for $\boldsymbol{\omega} \in \R^3$
with $\theta = \norm{\boldsymbol{\omega}}$,
\[
  \exp([\boldsymbol{\omega}]_\times) = I + \frac{\sin\theta}{\theta}
  [\boldsymbol{\omega}]_\times + \frac{1-\cos\theta}{\theta^2}
  [\boldsymbol{\omega}]_\times^2,
\]
where $[\boldsymbol{\omega}]_\times$ is the antisymmetric matrix of
$\boldsymbol{\omega}$. The formula is efficient and free of the singularities of
Euler angles, making it central to $\SE(3)$-equivariant networks on point
clouds.
\end{example}

\begin{figure}[h]
\centering
\begin{tikzpicture}[>=Stealth, scale=1.0]
\begin{scope}[shift={(-3.2,0)}]
  \draw[gdlgray!40] (0,0) circle (1.3cm);
  \draw[->] (-1.7,0) -- (1.7,0) node[right, font=\tiny] {$x$};
  \draw[->] (0,-1.7) -- (0,1.7) node[above, font=\tiny] {$y$};
  \draw[->, very thick, gdlblue!70] (0,0) -- (1.3,0);
  \draw[->, very thick, gdlred!70] (0,0) -- ({1.3*cos(50)},{1.3*sin(50)});
  \draw[gdlpurple!70] (0.7,0) arc (0:50:0.7) ;
  \node[gdlpurple!70, font=\tiny] at (0.95,0.35) {$\theta$};
  \node[font=\small\bfseries] at (0,-2.1) {$\SO(2)$};
\end{scope}
\begin{scope}[shift={(2.6,0)}]
  \draw[->] (-1.8,0) -- (1.8,0) node[right, font=\tiny] {$x$};
  \draw[->] (0,-1.6) -- (0,1.9) node[above, font=\tiny] {$z$};
  \draw[->] (1.0,0.6) -- (-1.1,-0.7) node[left, font=\tiny] {$y$};
  \shade[ball color=gdlorange!20, opacity=0.6] (0,0) circle (1.3cm);
  \draw[->, very thick, gdlblue!70] (0,0) -- (1.2,0.55);
  \draw[->, thick, gdlgreen!60!black] (0,0) -- (0,1.5) node[above, font=\tiny] {$R_z$};
  \draw[->, gdlpurple!70] (0.55,0.25) arc (24:80:0.6);
  \node[font=\small\bfseries] at (0,-2.1) {$\SO(3)$};
\end{scope}
\end{tikzpicture}
\caption[Rotations in the plane and in space]{Left: the group $\SO(2)$ of planar
rotations $R_\theta$, topologically the circle $S^1$. Right: the group $\SO(3)$
of spatial rotations; the elementary matrices $R_x, R_y, R_z$ rotate about the
coordinate axes.}
\label{fig:found:rotations}
\end{figure}

\begin{definition}[Euclidean groups]
\label{def:found:euclidean-groups}
The \emph{special Euclidean group} $\SE(n) = \R^n \rtimes \SO(n)$ of rigid
motions (rotation and translation) and the \emph{Euclidean group}
$E(n) = \R^n \rtimes \Orth(n)$ (also including reflections) are Lie groups of
dimension $\frac{n(n+1)}{2}$.
\end{definition}

Elements of $\SE(3)$ are conveniently written in homogeneous coordinates as
$\left(\begin{smallmatrix} R & t \\ 0 & 1\end{smallmatrix}\right)$ with
$R \in \SO(3)$, $t \in \R^3$, acting on $x$ by $x \mapsto Rx + t$. The Lie algebra
$\Lie{se}(3)$ has dimension $6$, with three rotation and three translation
generators. In geometric deep learning $\SE(3)$ is essential for 3D point clouds
and molecular data, where predictions must be invariant to rigid motions.

\begin{definition}[Linear and unitary groups]
\label{def:found:linear-unitary}
The \emph{general linear group} $\GL(n,\R) = \{A : \det A \ne 0\}$ is a
non-compact Lie group of dimension $n^2$; the \emph{special linear group}
$\SL(n,\R) = \{A : \det A = 1\}$ has dimension $n^2 - 1$. The \emph{unitary group}
$\U(n) = \{U \in \C^{n\times n} : U^\dagger U = I_n\}$ is compact of dimension
$n^2$, and the \emph{special unitary group} $\SU(n) = \{U \in \U(n) : \det U =
1\}$ has dimension $n^2 - 1$.
\end{definition}

The group $\GL(n,\R)$ is an open subset of $\R^{n^2}$, hence of dimension $n^2$;
for $\SL(n,\R)$ the regular level set theorem with $F = \det$ gives $n^2 - 1$. A
fundamental fact is the relation between $\SU(2)$ and $\SO(3)$: the adjoint action
$\Ad_U(X) = UXU^\dagger$ of $\SU(2) \cong S^3$ on its Lie algebra $\Lie{su}(2)
\cong \R^3$ defines a surjective homomorphism $\pi : \SU(2) \to \SO(3)$ with
kernel $\{I, -I\} \cong \Z_2$, so by the first isomorphism theorem
$\SO(3) \cong \SU(2)/\Z_2$: $\SU(2)$ is the double cover of $\SO(3)$, and
$\Lie{su}(2) \cong \Lie{so}(3)$ as Lie algebras. This is the reason unit
quaternions ($\cong \SU(2)$) give a singularity-free parametrization of 3D
rotations.

\section{Representations}
\label{sec:found:representations}

A group action moves the points of a domain around. Inside a neural network,
however, symmetries must act on \emph{feature spaces} --- vector spaces of
activations. The notion that makes a symmetry act linearly is a representation.

\begin{definition}[Representation]
\label{def:found:representation}
A \emph{representation} of a group $\group$ on a vector space $V$ over a field
$\mathbb{K}$ is a homomorphism $\rho : \group \to \GL(V)$, where $\GL(V)$ is the
group of invertible linear maps on $V$. Equivalently, each $g$ is assigned an
invertible matrix $\rho(g)$ with $\rho(g_1 g_2) = \rho(g_1)\rho(g_2)$. For a Lie
group we additionally require $\rho$ to be smooth.
\end{definition}

Representations translate abstract symmetries into concrete linear operations on
the channels of a network.

\begin{proposition}[Induced Lie-algebra representation]
\label{prop:found:drho}
A smooth representation $\rho : G \to \GL(V)$ of a Lie group induces a
representation of its Lie algebra, $d\rho : \Lie{g} \to \Lie{gl}(V)$ with
$d\rho(X) = \frac{d}{dt}\big|_0\, \rho(\exp(tX))$, which preserves the bracket:
$d\rho([X,Y]) = [d\rho(X), d\rho(Y)]$.
\end{proposition}

\begin{proof}
Since $\rho$ and $\exp$ are smooth, $t \mapsto \rho(\exp(tX))$ is a smooth curve
in $\GL(V)$ and its derivative at $0$ lies in $\Lie{gl}(V) = T_{\mathrm{Id}}\GL(V)$.
Linearity in $X$ follows from the Baker--Campbell--Hausdorff formula $\exp(t(X+Y))
= \exp(tX)\exp(tY)\exp(O(t^2))$ on differentiating at $0$. For the bracket,
$\exp(tX)\exp(tY)\exp(-tX)\exp(-tY) = \exp(t^2[X,Y] + O(t^3))$; applying the
homomorphism $\rho$ and differentiating twice yields
$d\rho([X,Y]) = d\rho(X)d\rho(Y) - d\rho(Y)d\rho(X) = [d\rho(X), d\rho(Y)]$. See
\citep[Thm.~3.28]{Kirillov2008}.
\end{proof}

This is what lets us optimize in the linear algebra and act in the group. The
basic question of representation theory is how a representation decomposes.

\begin{definition}[Irreducible representation]
\label{def:found:irreducible}
A representation $\rho : \group \to \GL(V)$ with $V \ne \{0\}$ is
\emph{irreducible} if the only subspaces $W \subseteq V$ with
$\rho(g)W \subseteq W$ for all $g$ are $\{0\}$ and $V$.
\end{definition}

Irreducible representations are the building blocks from which all others are
assembled. For $\SO(2)$ the complex irreducibles are one-dimensional, indexed by
$n \in \Z$:
\[
  \rho_n : \SO(2) \to \C^*, \qquad \rho_n(R_\theta) = e^{i n\theta}.
\]
Each $n$ is a \emph{frequency}, and these are exactly the basis functions of the
classical Fourier series.

\begin{definition}[Compact group and Haar measure]
\label{def:found:haar}
A Lie group $G$ is \emph{compact} if it is compact as a topological space (e.g.\
$\SO(n)$, $\Orth(n)$, $\U(n)$, $\SU(n)$, and all finite groups). A compact group
carries a unique invariant probability measure, the \emph{Haar measure} $\mu$,
satisfying $\mu(gA) = \mu(Ag) = \mu(A)$ and $\mu(\group) = 1$.
\end{definition}

The Haar measure provides a canonical way to ``average over all symmetries.'' In
geometric deep learning this averaging is exactly what underlies group
convolution: the convolution integrates a filter against the signal using the
Haar measure, and its invariance is what guarantees equivariance of the resulting
operator~\citep{cohen2016group}.

\begin{theorem}[Complete reducibility]
\label{thm:found:decomposition}
Every finite-dimensional representation of a compact group $\group$ decomposes
as a direct sum of irreducibles, $V = \bigoplus_i n_i V_i$, where each $V_i$ is
irreducible and $n_i \in \mathbb{N}$ is its multiplicity.
\end{theorem}

\begin{proof}
Given any inner product $\inner{\cdot}{\cdot}_0$ on $V$, the average
$\inner{u}{v} = \int_\group \inner{\rho(g)u}{\rho(g)v}_0\, d\mu(g)$ is a
$\group$-invariant inner product: positivity holds since $\rho(g)$ is invertible,
and invariance follows from the change of variable $g' = gh$ and the invariance
of $\mu$. With respect to it, if $W$ is an invariant subspace then so is $W^\perp$
(for $v \in W^\perp$, $g \in \group$, $w \in W$: $\inner{\rho(g)v}{w} =
\inner{v}{\rho(g^{-1})w} = 0$ since $\rho(g^{-1})w \in W$). Induction on
$\dim V$: if $\rho$ is irreducible we are done; otherwise an invariant proper
$W$ splits off $V = W \oplus W^\perp$ with both summands of smaller dimension,
which by induction decompose.
\end{proof}

The maps between irreducibles are severely constrained, a fact known as Schur's
lemma.

\begin{lemma}[Schur]
\label{lem:found:schur}
Let $(\rho, V)$ and $(\sigma, W)$ be irreducible and let $T : V \to W$ satisfy
$T\rho(g) = \sigma(g)T$ for all $g$ (an intertwiner). Then $T$ is either zero or
an isomorphism. If moreover $V = W$ and the field is $\C$, then $T =
\lambda\,\mathrm{Id}$ for some scalar $\lambda$.
\end{lemma}

\begin{proof}
The kernel of $T$ is invariant ($v \in \ker T \Rightarrow T\rho(g)v =
\sigma(g)Tv = 0$), so by irreducibility of $\rho$ it is $\{0\}$ or $V$; the image
is invariant, so by irreducibility of $\sigma$ it is $\{0\}$ or $W$. Hence
$T = 0$ or $T$ is a bijection. Over $\C$, $T$ has an eigenvalue $\lambda$; then
$T - \lambda\,\mathrm{Id}$ is an intertwiner with nontrivial kernel, so it
vanishes.
\end{proof}

Schur's lemma is the algebraic root of the parameter efficiency of equivariant
networks: an equivariant linear layer between representations is forced to be
block-diagonal in the irreducible decomposition, with each block proportional to
the identity, so the number of free parameters is dictated by the symmetry rather
than by the input size. The \emph{matrix coefficients} $\rho_{ij}(g) =
\inner{e_i}{\rho(g)e_j}$ play the role of coordinate functions; the Schur
orthogonality relations state that, for irreducible unitary representations,
$\int_G \overline{\rho_{ij}(g)}\sigma_{kl}(g)\, d\mu(g) = \delta_{\rho\sigma}
\delta_{ik}\delta_{jl}/d_\rho$, so the normalized matrix coefficients form an
orthonormal system in $L^2(G)$. This decomposition culminates in the
Peter--Weyl theorem.

\begin{theorem}[Peter--Weyl]
\label{thm:found:peterweyl}
Let $\group$ be a compact Lie group. The space $L^2(\group)$ decomposes as
\[
  L^2(\group) = \bigoplus_{\lambda \in \widehat{\group}} d_\lambda\, V_\lambda,
\]
where $\widehat{\group}$ is the set of equivalence classes of irreducible
representations, $V_\lambda$ is the space of $\lambda$, and $d_\lambda = \dim
V_\lambda$.
\end{theorem}

\begin{proof}[Sketch]
By the Schur orthogonality relations the normalized matrix coefficients
$\{\sqrt{d_\lambda}\,\rho^\lambda_{ij}\}$ form an orthonormal system in
$L^2(\group)$. Completeness follows because the representations of a compact
group separate points, so by the Stone--Weierstrass theorem the matrix
coefficients are dense in $C(\group)$, hence in $L^2(\group)$: any $f$ orthogonal
to all of them makes the operator $T_f = \int_\group f(g)\rho^\lambda(g)\,d\mu(g)$
vanish for every $\lambda$, forcing $f = 0$.
\end{proof}

The Peter--Weyl theorem is the generalization of the Fourier series to
non-abelian groups: the matrix entries of the irreducibles play the role of the
exponentials $e^{in\theta}$ and form a complete orthonormal basis of
$L^2(\group)$. For $\group = \SO(2)$ it reduces exactly to the classical Fourier
expansion $f(\theta) = \sum_n \hat f_n\,e^{in\theta}$. For network design it has
three consequences, developed in \Cref{ch:groups}: the channels of an
equivariant layer are indexed by irreducibles $\lambda$; the per-channel
parameter count $d_\lambda^2$ is fixed by the group, independent of input
resolution; and equivariant convolutions become pointwise multiplications in this
generalized Fourier (spectral) domain. Concretely, every $G$-equivariant linear
operator $\Phi$ on $L^2(G)$ preserves the isotypic decomposition and acts on each
block as $I_{d_\lambda} \otimes W_\lambda$ for a matrix $W_\lambda$, by Schur's
lemma.

\section{Homogeneous spaces and quotients}
\label{sec:found:homogeneous}

So far symmetries have acted on data. We now turn the construction around and use
a group to \emph{build} the domain itself. Many of the spaces on which geometric
data lives --- the sphere $S^2$, projective spaces, Grassmannians --- share a
hidden structure: they are quotients $G/H$ of a Lie group by a closed subgroup.
Recognizing a domain as such a quotient is the single most useful step in
designing an equivariant architecture for it.

\begin{definition}[Orbit and stabilizer]
\label{def:found:orbit}
Let $\group$ act on a set $X$ and let $x \in X$. The \emph{orbit} of $x$ is
$\Orb(x) = \{g \cdot x : g \in \group\}$, the points reachable from $x$. The
\emph{stabilizer} of $x$ is $\Stab(x) = \{g \in \group : g \cdot x = x\}$, the
subgroup fixing $x$.
\end{definition}

\begin{theorem}[Orbit--stabilizer]
\label{thm:found:orbitstab}
Let $\group$ act on $X$ and let $x \in X$. There is a natural bijection
\[
  \group / \Stab(x) \;\cong\; \Orb(x), \qquad g\,\Stab(x) \mapsto g \cdot x.
\]
\end{theorem}

\begin{proof}
The map $\psi(g\,\Stab(x)) = g\cdot x$ is well defined: if $g\,\Stab(x) =
h\,\Stab(x)$ then $h^{-1}g \in \Stab(x)$, so $g\cdot x = h\cdot((h^{-1}g)\cdot x)
= h\cdot x$. It is injective by the same computation read backwards, and
surjective because every point of $\Orb(x)$ is $g\cdot x$ for some $g$.
\end{proof}

\begin{figure}[h]
\centering
\begin{tikzpicture}[
    >=Stealth,
    pt/.style={circle, fill=#1, minimum size=7pt, inner sep=0pt}
]
\draw[thick, fill=gdlgreen!12, rounded corners] (-5.2,-1.6) rectangle (-3.4,2.1);
\node[above, font=\small\bfseries, text=gdlgreen!60!black] at (-4.3,2.1) {Group $\group$};
\foreach \y/\lab in {1.3/$g_1\Stab(x)$, 0.3/$g_2\Stab(x)$, -0.7/$\Stab(x)$} {
    \draw[thick, fill=gdlgreen!22, rounded corners=1pt] (-5.05,\y) rectangle (-3.55,\y-0.55);
    \node[font=\tiny] at (-4.3,\y-0.27) {\lab};
}
\draw[->, very thick, gdlred!75!black] (-3.2,0.4) -- node[above, font=\small] {$\cong$} (-1.2,0.4);
\draw[thick, fill=gdlgray!12, rounded corners] (-1,-1.6) rectangle (4.6,2.6);
\node[above, font=\small\bfseries, text=gdlgray!50!black] at (1.8,2.6) {Space $X$};
\draw[gdlblue!60, thick, fill=gdlblue!12] (1.7,0.5) ellipse (1.9cm and 0.95cm);
\node[above, font=\tiny, text=gdlblue!70!black] at (1.7,1.5) {$\Orb(x) = \group\cdot x$};
\node[pt=gdlred!75] (x) at (0.4,0.5) {};
\node[below left, font=\small] at (x) {$x$};
\node[pt=gdlorange!80] (gx1) at (2.2,1.1) {};
\node[above, font=\small] at (gx1) {$g_1\cdot x$};
\node[pt=gdlorange!80] (gx2) at (2.7,-0.1) {};
\node[below, font=\small] at (gx2) {$g_2\cdot x$};
\draw[->, thick, gdlpurple!75] (x) to[bend left=28] node[above, font=\tiny] {$g_1$} (gx1);
\draw[->, thick, gdlpurple!75] (x) to[bend right=18] node[below, font=\tiny] {$g_2$} (gx2);
\draw[->, thick, gdlgreen!55!black, loop above, looseness=6] (x) to node[above, font=\tiny] {$h\in\Stab(x)$} (x);
\end{tikzpicture}
\caption[Orbit--stabilizer theorem]{The orbit--stabilizer theorem. Each coset
$g\,\Stab(x)$ corresponds to exactly one point $g\cdot x$ of the orbit, giving the
bijection $\group/\Stab(x) \cong \Orb(x)$.}
\label{fig:found:orbit}
\end{figure}

When $\group$ is a Lie group acting transitively and $H$ is a closed subgroup,
$G/H$ is a smooth manifold, a \emph{homogeneous space}, on which $\group$ acts
transitively by $g'\cdot(gH) = (g'g)H$. The orbit--stabilizer theorem says that
any domain admitting a transitive group of symmetries is a homogeneous space:
identify the symmetry group $\group$, fix a reference point, and the stabilizer
$H$ of that point determines the domain as $G/H$.

\begin{example}[The sphere as $\SO(3)/\SO(2)$]
\label{ex:found:sphere}
Let $\SO(3)$ act on $S^2$ by rotation. The action is transitive: any unit vector
can be rotated to any other (complete each to a positively oriented orthonormal
basis and rotate one basis to the other). Fix the north pole $p = (0,0,1)^\top$;
a rotation fixes $p$ exactly when its axis is the $z$-axis, and these rotations
form a subgroup isomorphic to $\SO(2)$. Hence $\Stab(p) \cong \SO(2)$, and the
orbit--stabilizer theorem gives $S^2 \cong \SO(3)/\SO(2)$. More generally,
$S^{n-1} \cong \SO(n)/\SO(n-1)$. A signal $f : S^2 \to \R^C$ lifts to a function
on $\SO(3)$ that is right-invariant under $\SO(2)$, namely
$\tilde f(g) = f(g\cdot p)$, so the natural equivariant convolution on the sphere
is a correlation over the full group $\SO(3)$ rather than over a flat projection
--- the construction underlying spherical networks of
\Cref{ch:geodesics,ch:applications}.
\end{example}

This example illustrates a general recipe for geometric architectures: (1)
identify the data domain $\Omega$; (2) determine the relevant symmetry group
$\group$; (3) compute the stabilizer $H$ of a reference point; (4) recognize
$\Omega \cong G/H$ by the orbit--stabilizer theorem; and (5) build the
equivariant architecture using the harmonic analysis of $\group$
(\Cref{thm:found:peterweyl}). The same recipe identifies $\R^n \cong E(n)/\Orth(n)
\cong \SE(n)/\SO(n)$ (the base case of classical convolutional networks), the
real projective space $\mathbb{RP}^{n-1} \cong \Orth(n)/(\Orth(n-1) \times \Z_2)$,
and the Grassmannian $Gr(k,n) \cong \Orth(n)/(\Orth(k) \times \Orth(n-k))$.

\begin{proposition}[Geometric properties of homogeneous spaces]
\label{prop:found:homogeneous-props}
Let $G/H$ be a homogeneous space. Then $G$ acts transitively on $G/H$; if $G$ is
connected then $G/H$ is connected; and $G$ acts by isometries of any
$G$-invariant Riemannian metric on $G/H$ (a metric $g$ with $g_p(u,v) =
g_{\gamma\cdot p}(d\gamma_p u, d\gamma_p v)$ for all $\gamma \in G$).
\end{proposition}

\begin{proof}
Transitivity: given $g_1 H, g_2 H$, take $g = g_2 g_1^{-1}$, so $g\cdot(g_1 H) =
g_2 H$. Connectedness: the canonical projection $\pi : G \to G/H$ is continuous
and surjective, so $G/H = \pi(G)$ is connected when $G$ is. Invariance of metric:
$G$-invariance says exactly that each translation $L_\gamma : p \mapsto
\gamma\cdot p$ satisfies $L_\gamma^* g = g$, i.e.\ is an isometry.
\end{proof}

On a homogeneous space an equivariant operator is completely determined by its
behaviour at the base point $o = eH$: since $\group$ acts transitively, the value
of $\Phi f$ at any $x = g\cdot o$ is fixed by equivariance,
$(\Phi f)(x) = \rho_2(g)(\Phi(L_{g^{-1}}f))(o)$. A short computation shows that
the only equivariant linear operators on $L^2(G)$ are generalized convolutions:
the kernel must satisfy $k(hg, hg') = k(g, g')$, so $k(g,g') = \kappa(g^{-1}g')$
and $(\Phi f)(g) = \int_G \kappa(g^{-1}g')f(g')\,d\mu(g') = (\kappa * f)(g)$. This
is the mathematical reason convolution is the canonical operation on symmetric
domains, a thread we follow through \Cref{ch:grids,ch:groups,ch:graphs}. The
choice of homogeneous space thus determines the symmetry group $G$ and stabilizer
$H$, which in turn dictate the structure of the available equivariant
representations.

\medskip

The three threads of this chapter --- groups, representations, and manifolds ---
meet in the Lie group: a manifold whose tangent space at the identity is its Lie
algebra, with the exponential map as the canonical curve through the identity in
each direction. Together with the neural-network foundations of the first half,
they furnish the complete vocabulary for the geometric priors of \Cref{ch:priors}
and the architectures that follow.
